\documentclass[a4paper,12pt]{article}

\usepackage[
  a4paper,
  left=2.1cm,
  right=2.1cm,
  top=2.3cm,
  bottom=2.3cm
]{geometry}
\usepackage{setspace}
\usepackage{microtype}
\usepackage{amsmath,amssymb,amsfonts,amsthm,mathtools}
\usepackage{bm}
\usepackage{mathrsfs}
\usepackage{graphicx}
\usepackage{float}
\usepackage{enumitem}
\usepackage[dvipsnames]{xcolor}
\usepackage[T1]{fontenc}
\usepackage[utf8]{inputenc}
\usepackage[english]{babel}
\usepackage[round]{natbib}
\usepackage[colorlinks=true, linkcolor=MidnightBlue, citecolor=MidnightBlue, urlcolor=blue]{hyperref}
\usepackage{soul}
\usepackage{subcaption}

\setstretch{1.1}

\newtheorem{theorem}{Theorem}[section]
\newtheorem{proposition}[theorem]{Proposition}
\newtheorem{lemma}[theorem]{Lemma}
\newtheorem{corollary}[theorem]{Corollary}
\newtheorem{assumption}[theorem]{Assumption}
\theoremstyle{definition}
\newtheorem{definition}[theorem]{Definition}
\newtheorem{example}[theorem]{Example}
\theoremstyle{remark}
\newtheorem{remark}[theorem]{Remark}

\DeclareMathOperator{\VaR}{VaR}
\DeclareMathOperator{\TVaR}{TVaR}
\DeclareMathOperator{\Cov}{Cov}

\newcommand{\R}{\mathbb{R}}
\newcommand{\E}{\mathbb{E}}
\newcommand{\Var}{\mathrm{Var}}
\newcommand{\1}{\mathbf{1}}

\DeclareMathOperator*{\argmin}{arg\,min}
\DeclareMathOperator*{\argmax}{arg\,max}


\title{Voluntary Partial Participation in Peer-to-Peer Insurance with Heterogeneous Risks}
\date{\today}

\begin{document}

\maketitle

\begin{abstract}
We study voluntary participation in peer-to-peer (P2P) insurance pools in which 
agents differ in their risk characteristics. A finite population is composed of 
high-risk and low-risk participants who choose type-contingent contribution 
rates \((\alpha_H,\alpha_L)\) in a common budget-balanced sharing rule. Each 
type's participation is constrained by individual rationality relative to 
autarky and, when contributions are interpreted as a type-contingent contract 
menu, by incentive compatibility. Within this framework, we ask which 
heterogeneous pools can sustain a P2P arrangement, how the sharing rule and 
pool composition jointly determine welfare, and how decentralized participation 
differs from a planner allocation.

We obtain three sets of results. First, we develop a general structural 
framework: along any participation direction, welfare admits a quadratic 
representation, is concave, and---under type-neutral welfare weights---depends 
only on variance components, because budget balance neutralizes the mean. 
Second, in the independent-risk benchmark, the low-risk participation 
constraint generates a tractable frontier in the \((\alpha_H,\alpha_L)\) plane: 
the feasible set is star-shaped, the planner's problem reduces to boundary 
optimization, and the feasibility of full pooling is governed by an explicit 
threshold that becomes asymptotically tight in large populations. Third, with 
correlated risks, covariance exposure between an agent's own loss and the pool 
becomes a separate determinant of feasibility and welfare, and may either 
contract or expand the feasible set depending on the sign of cross-type 
dependence.

Comparing planner and decentralized participation reveals a composition 
externality: each type internalizes its own insurance gain and redistribution 
cost, but not the effect of its participation on the other type's risk 
exposure. In the independent benchmark, this externality has a definite sign 
and produces low-risk underparticipation; under dependence its sign is 
covariance-contingent. The economic implication is that the central design 
challenge in P2P insurance is not pool size but pool composition: practical 
arrangements must be engineered around the low-risk participation constraint 
that ultimately determines whether diversification gains can be realized.
\end{abstract}

\noindent\textbf{Keywords:} peer-to-peer insurance; risk sharing; participation; heterogeneous risks.\\


\section{Introduction}

Peer-to-peer insurance is attractive because it replaces part of conventional insurance with mutual risk sharing inside a pool. The difficulty is that the pool is rarely homogeneous. Participants may differ in expected losses, loss variances, and dependence with the rest of the group. These differences matter not only for how losses should be allocated after they occur, but also for whether different risk categories are willing to participate in the first place.

This paper studies this participation problem in a finite heterogeneous P2P insurance pool. Pooling can reduce idiosyncratic risk, but it can also redistribute expected losses across types. A low-risk participant may gain from diversification while at the same time being exposed to claims generated by high-risk participants. A high-risk participant may benefit from broad participation, but her participation changes the terms faced by the rest of the pool. Hence the design problem combines choice of a sharing rule with the determination of which type-contingent participation profiles can be sustained.

We consider a two-type model with high-risk and low-risk agents. Agents have mean--variance preferences and choose, or are assigned, type-contingent participation rates in a common budget-balanced pool. A type-\(k\) agent contributes a fraction \(\alpha_k\) of her loss to the pool and receives a share of the pooled loss determined by the same contribution weights. The pair \((\alpha_H,\alpha_L)\) therefore describes both the intensity of pooling and the composition of participation. Feasible allocations must satisfy participation constraints, and, when participation rates are interpreted as type-contingent contract terms, the relevant incentive constraints as well.

The paper is related to the classical literature on optimal insurance and risk sharing, beginning with \citet{Borch1962} and \citet{Arrow1974}, and to subsequent work on efficient sharing rules and insurance design under alternative constraints \citep{Raviv1979,CumminsMahul2004,DanaScarsini2007,JouiniSchachermayerTouzi2008}. That literature provides the efficiency benchmark for allocating risk, but it usually abstracts from the finite-pool participation problem that arises when heterogeneous agents voluntarily enter a common P2P arrangement.

The analysis also connects to recent work on P2P insurance, mutual aid, and decentralized risk sharing. One line of work studies conditional-mean and related sharing rules and their fairness or stochastic-order properties \citep{DenuitDhaene2012,Denuit2019,DenuitEtAl2023}. Another line studies P2P insurance with outside coverage, deposits, mutual aid, or decentralized pooling mechanisms \citep{abdikerimova2022peer,FengEtAl2023,LevantesiPiscopo2022,ChenFengMaoHu2024,BoonenChenHu2025}. The closest papers in this group ask how a pool should share losses, choose coverage, or allocate fixed and variable costs once the relevant participants are in the arrangement. The present paper asks a complementary question: how do heterogeneity and participation constraints restrict the set of pooling arrangements that can be sustained?

A third connection is with the literature on informal insurance, coalition stability, and limited commitment \citep{coate1993reciprocity,genicot2003group}. Those models emphasize that risk sharing may be valuable but difficult to sustain when agents can opt out. The mechanism in this paper is different: participation is constrained by type-specific risk characteristics and by the redistribution embedded in the P2P sharing rule. The common theme is that welfare-improving risk sharing need not be implementable once participation is voluntary.

The first part of the paper develops general structural results. Along any fixed participation direction, the sharing operator is linear in the scale of participation. As a result, planner welfare is quadratic along rays. Under type-neutral welfare weights, the mean component of aggregate welfare is invariant because the sharing rule is budget balanced; changes in welfare come through the variance terms. These observations identify the objects governing whether a welfare-improving allocation is feasible.

Under independent risks, the participation constraints take an explicit form. In that case, the participation constraints can be written explicitly, and the low-risk constraint generates a frontier in the \((\alpha_H,\alpha_L)\) plane. This frontier captures the main economic force in the model: broader pooling is valuable because it reduces idiosyncratic variance, but it is limited by the expected redistribution borne by low-risk agents. Full pooling is feasible only when the diversification gain is large enough to compensate the types whose participation constraints are most exposed to redistribution. In finite pools, this condition is substantive; it is not implied merely by the desirability of pooling from an aggregate point of view.

The paper then considers correlated losses. Dependence introduces a new object: the covariance between an individual's own loss and the aggregate pool. With independent risks, participation mainly changes retained variance and pooled variance. With correlated risks, it also changes covariance exposure. This additional term can shrink or expand the feasible set, alter the sign of welfare derivatives, and break the simple geometry obtained in the independent benchmark.

Finally, we compare the planner allocation with decentralized participation. In the planner problem, \(\alpha_H\) and \(\alpha_L\) are chosen to maximize aggregate welfare subject to feasibility. In the decentralized problem, each type chooses its own participation rate, taking the other type's rate as given. The difference between the two first-order conditions identifies the participation wedge: a type internalizes its own insurance gain and redistribution cost, but not the full effect of its participation on the other type's risk exposure. This wedge can lead to too little participation, too much participation, or distorted pool composition, depending on the covariance and redistribution effects.

The remainder of the paper is organized as follows. Section~\ref{sec:model} introduces the model and the participation constraints. Section~\ref{sec:general-results} establishes the general ray and welfare structure. Section~\ref{sec:ind} studies the independent-risk benchmark and characterizes the feasible frontier and planner allocation. Section~\ref{sec:correlated} extends the analysis to correlated risks. Section~\ref{sec:decentralized} compares planner and decentralized participation. Section~\ref{sec:numerical} provides numerical illustrations, and Section~\ref{sec:conclusion} concludes.
\section{General Setting}\label{sec:model}

\subsection{Model Setup}

There are two types of agents indexed by \(k\in\{H,L\}\), where \(H\) denotes high-risk and \(L\) denotes low-risk agents. Let \(n_k\in\mathbb{N}_+\) denote the number of type-\(k\) agents, and let \(n=n_H+n_L\) be the total population size. Agents are indexed by \(i\in\{1,\dots,n\}\), and \(\tau(i)\in\{H,L\}\) denotes the type of agent \(i\).

Each agent \(i\) faces a random loss \(X_i\). Let
$\mathbf X=(X_1,\dots,X_n)^\top$
denote the vector of losses. We assume that \(\mathbf X\) has finite second moments, with mean vector
$\boldsymbol{\mu}:=\E[\mathbf X]$
and covariance matrix
$\Sigma_X:=\Var(\mathbf X)$,
where \(\Sigma_X\) is symmetric and positive semidefinite.

Losses are type-dependent up to the second moment. Formally, there exist constants $\mu_k\in\mathbb{R}$ and $\sigma_k^2\ge 0$ for each $k\in\{H,L\}$, and constants $\sigma_{k\ell}\in\mathbb{R}$ for each $k,\ell\in\{H,L\}$, such that for every agent $i$,
\[
\E[X_i]=\mu_{\tau(i)},
\qquad
\Var(X_i)=\sigma_{\tau(i)}^2,
\]
and for every pair of distinct agents $i\neq j$,
\[
\Cov(X_i,X_j)=\sigma_{\tau(i)\tau(j)}.
\]
In other words, the first and second moments of losses depend only on agents' types.

We further assume that \(\mu_H>\mu_L>0\).

Agents have mean--variance preferences. For any random loss \(Y\),
\begin{equation}\label{eq:utility-general}
U(Y)=-\E[Y]-\frac{\gamma}{2}\Var(Y),
\end{equation}
where \(\gamma>0\) is the common coefficient of risk aversion. Initial wealth is normalized to zero.

Each type \(k\) chooses a contribution rate \(\alpha_k\in[0,1]\). An agent of type \(k\) contributes \(\alpha_k X_i\) to a common pool, retains \((1-\alpha_k)X_i\), and receives a share of the pool proportional to her contribution.

Let \(T:=\alpha_H n_H+\alpha_L n_L>0\), the realized net loss of agent \(i\) of type \(k\) is
\begin{equation}\label{eq:netloss-general}
\tilde X_i
=
(1-\alpha_k)X_i
+
\frac{\alpha_k}{T}\sum_{j=1}^n \alpha_{\tau(j)}X_j.
\end{equation}

In matrix form, let \(\boldsymbol{\alpha}\in[0,1]^n\) collect individual contribution rates, and let \(D_\alpha=\mathrm{diag}(\boldsymbol{\alpha})\). Then
\[
\tilde{\mathbf X}
=
A(\boldsymbol{\alpha})\mathbf X,
\qquad
A(\boldsymbol{\alpha})
=
I-D_\alpha+\frac{1}{T}\boldsymbol{\alpha}\boldsymbol{\alpha}^\top.
\]

\subsection{Participation, deviations, and planner’s problem}

Because agents are ex ante identical within type and choose common contribution rates, we focus on representative type-level objects. Define \(P:=\alpha_H n_H\mu_H+\alpha_L n_L\mu_L\).

Then a representative type-\(k\) agent has expected net loss
\begin{equation}\label{eq:mean-type-general}
\E[\tilde X_k]
=
(1-\alpha_k)\mu_k+\alpha_k\frac{P}{T}.
\end{equation}

To describe type-level variances under a structured covariance environment, we will later specialize to the type-symmetric covariance case. For now, we keep the covariance structure general.

We model participation choices as a simultaneous-move game in which each agent independently chooses her contribution rate. Let $S_{-i}:=\sum_{j\neq i}\alpha_j$ be the aggregate contribution weight of all agents other than agent $i$.

If a type-\(k\) agent deviates to an alternative contribution \(\alpha\in[0,1]\), her realized loss becomes
\begin{equation}\label{eq:devloss-general}
\tilde X_k(\alpha\mid \alpha_H,\alpha_L)
=
(1-\alpha)X_k
+
\frac{\alpha}{\alpha+S_{-i}}
\left(
\alpha X_k+\sum_{j\neq i}\alpha_j X_j
\right).
\end{equation}
Her deviation utility is
\begin{equation}\label{eq:devutility-general}
U_k(\alpha\mid \alpha_H,\alpha_L)
=
-\E[\tilde X_k(\alpha\mid \alpha_H,\alpha_L)]
-\frac{\gamma}{2}\Var(\tilde X_k(\alpha\mid \alpha_H,\alpha_L)).
\end{equation}


The planner chooses \((\alpha_H,\alpha_L)\in[0,1]^2\) to maximize weighted utilitarian welfare
\begin{equation}\label{eq:welfare-general}
W(\alpha_H,\alpha_L)
=
n_H\omega_H U_H
+
n_L\omega_L U_L,
\end{equation}
where \(\omega_H,\omega_L>0\) are welfare weights and \(U_k\) denotes truthful utility.


Feasible allocations must satisfy individual rationality (IR) and incentive compatibility (IC). 
We interpret \((\alpha_H,\alpha_L)\) as a type-contingent contribution menu: a type-\(k\) agent is assigned the contribution rate \(\alpha_k\), taking the contribution choices of all other agents as given.

The outside option is autarky, corresponding to zero participation in the pool. 
Thus, if an agent deviates to \(\alpha=0\), she retains her own loss and makes no transfer to or from the pool, while all other agents continue to contribute according to the candidate allocation \((\alpha_H,\alpha_L)\).

\paragraph{Individual rationality (IR).}
A type-\(k\) agent weakly prefers participating truthfully to autarky:
\begin{equation}\label{eq:IR-k}
U_k(\alpha_k \mid \alpha_H,\alpha_L)\ge U_k(0 \mid \alpha_H,\alpha_L),
\qquad k\in\{H,L\}.
\end{equation}
Equivalently, in type-specific form,
\begin{align}
U_H(\alpha_H \mid \alpha_H,\alpha_L) &\ge U_H(0 \mid \alpha_H,\alpha_L), \label{eq:IR-H}\\
U_L(\alpha_L \mid \alpha_H,\alpha_L) &\ge U_L(0 \mid \alpha_H,\alpha_L). \label{eq:IR-L}
\end{align}

\paragraph{Incentive compatibility (IC).}
Given the menu \((\alpha_H,\alpha_L)\), each type weakly prefers its own prescribed contribution rate to mimicking the other type's rate. Hence,
\begin{subequations}\label{eq:IC}
\begin{align}
U_H(\alpha_H \mid \alpha_H,\alpha_L) 
&\ge U_H(\alpha_L \mid \alpha_H,\alpha_L), \label{eq:IC-H}\\
U_L(\alpha_L \mid \alpha_H,\alpha_L) 
&\ge U_L(\alpha_H \mid \alpha_H,\alpha_L). \label{eq:IC-L}
\end{align}
\end{subequations}

That is, the high-risk type does not want to imitate the low-risk contribution rule, and the low-risk type does not want to imitate the high-risk contribution rule.

\section{Structural Results in the General Framework}\label{sec:general-results}

We first record several structural properties that do not depend on a specific covariance structure. These results do not fully characterize feasibility under general dependence, but they isolate the geometric and welfare objects used in the sharper benchmark analysis below.

\subsection{Welfare along rays}

Fix a direction \(\bar{\boldsymbol{\alpha}}\) and consider the ray
\[
\boldsymbol{\alpha}(t)=t\bar{\boldsymbol{\alpha}},
\qquad t\in[0,1].
\]
Because the sharing rule is homogeneous in contribution rates, the induced sharing operator is linear in \(t\) along any such ray.

Set \(\bar T:=\1^\top\bar{\boldsymbol{\alpha}}>0\) and define
\(B(\bar{\boldsymbol{\alpha}})
=
- D_{\bar{\boldsymbol{\alpha}}}
+
\frac{1}{\bar T}\bar{\boldsymbol{\alpha}}\bar{\boldsymbol{\alpha}}^\top\).
Then, since \(T(t)=\1^\top\boldsymbol{\alpha}(t)=t\bar T\), we have
\(A(\boldsymbol{\alpha}(t))=I+tB(\bar{\boldsymbol{\alpha}})\).
It follows that welfare is quadratic in the ray parameter.

\begin{proposition}\label{prop:ray-welfare-general}
Fix a direction \(\bar{\boldsymbol{\alpha}}\) with \(\bar T>0\). Along the ray \(t\mapsto t\bar{\boldsymbol{\alpha}}\), the planner welfare takes the form
\[
W(t)=W(0)+t\,\ell(\bar{\boldsymbol{\alpha}})-\frac{\gamma}{2}t^2 q(\bar{\boldsymbol{\alpha}}),
\]
where
\begin{align*}
\ell(\bar{\boldsymbol{\alpha}})
&=
-\mathbf m^\top B(\bar{\boldsymbol{\alpha}})\boldsymbol{\mu}
-
\gamma\operatorname{tr}\!\big(\Omega B(\bar{\boldsymbol{\alpha}})\Sigma_X\big),\\
q(\bar{\boldsymbol{\alpha}})
&=
\operatorname{tr}\!\big(\Omega B(\bar{\boldsymbol{\alpha}})\Sigma_X B(\bar{\boldsymbol{\alpha}})^\top\big),
\end{align*}
with \(\Omega\) the diagonal matrix of welfare weights replicated at the individual level, and \(\mathbf m\) the corresponding vector of welfare weights.
\end{proposition}

\begin{proof}
Let \(B:=B(\bar{\boldsymbol{\alpha}})\). Since \(A(\boldsymbol{\alpha}(t))=I+tB\), we have
\(\tilde{\mathbf X}(t)=(I+tB)\mathbf X\), and hence
\(\E[\tilde{\mathbf X}(t)]=(I+tB)\boldsymbol{\mu}\).
and
\[
\Var(\tilde{\mathbf X}(t))
=
(I+tB)\Sigma_X(I+tB)^\top
=
\Sigma_X+t(B\Sigma_X+\Sigma_XB^\top)+t^2B\Sigma_XB^\top.
\]
Substituting these expressions into the welfare functional and collecting terms in \(t\) and \(t^2\) yields the stated formula.
\end{proof}

The next lemma shows that welfare is concave along every ray.
\begin{lemma}\label{lem:concavity-general}
Suppose that the welfare weights are nonnegative. Then for every ray direction
\(\bar{\boldsymbol{\alpha}}\),
\[
q(\bar{\boldsymbol{\alpha}})
=
\operatorname{tr}\!\bigl(\Omega B(\bar{\boldsymbol{\alpha}})\Sigma_X
B(\bar{\boldsymbol{\alpha}})^\top\bigr)
\ge 0.
\]
The welfare is concave along every ray and strictly concave whenever
\[
\Omega^{1/2}B(\bar{\boldsymbol{\alpha}})\Sigma_X^{1/2}\neq 0.
\]
\end{lemma}

\begin{proof}
Let \(B:=B(\bar{\boldsymbol{\alpha}})\). Since \(\Omega\) is diagonal with
nonnegative entries, it is positive semidefinite; \(\Sigma_X\) is positive
semidefinite as a covariance matrix. Therefore,
\begin{align*}
q(\bar{\boldsymbol{\alpha}})
&=
\operatorname{tr}\!\bigl(\Omega B\Sigma_X B^\top\bigr) \\
&=
\operatorname{tr}\!\bigl(
(\Omega^{1/2}B\Sigma_X^{1/2})
(\Omega^{1/2}B\Sigma_X^{1/2})^\top
\bigr) \\
&=
\bigl\|\Omega^{1/2}B\Sigma_X^{1/2}\bigr\|_F^2
\ge 0.
\end{align*}
By Proposition~\ref{prop:ray-welfare-general},
\[
W(t)=W(0)+t\,\ell(\bar{\boldsymbol{\alpha}})
-\frac{\gamma}{2}t^2 q(\bar{\boldsymbol{\alpha}}),
\]
so \(W\) is concave in \(t\), and strictly concave if
\(q(\bar{\boldsymbol{\alpha}})>0\).
\end{proof}

\subsection{A conditional boundary principle}

The quadratic ray representation immediately yields a conditional boundary-optimality result.

\begin{proposition}\label{prop:boundary-general}
Let \(\mathcal F\subseteq[0,1]^2\) denote the feasible set. Suppose:
\begin{enumerate}[label=(\roman*), leftmargin=2em]
    \item \(\mathcal F\) is star-shaped with respect to the origin;
    \item welfare is concave along each feasible ray;
    \item along each feasible ray, the welfare-maximizing point weakly lies at or beyond the boundary of the feasible segment.
\end{enumerate}
Then every welfare-maximizing feasible allocation lies on \(\partial \mathcal F\).
\end{proposition}

\begin{proof}
By star-shapedness, every feasible point lies on a feasible ray from the origin to a boundary point. By assumption (iii), welfare is weakly increasing along each such feasible ray segment. Hence every interior feasible point is weakly dominated by a boundary point on the same ray. Therefore, an optimal allocation must be on \(\partial\mathcal F\).
\end{proof}


The conditions of Proposition~\ref{prop:boundary-general} are verified directly in the independent benchmark (Section~\ref{sec:ind}); additional restrictions on the covariance structure are required in the correlated case.

\subsection{Mean invariance under type-neutral welfare weights}

Under type-neutral welfare weights, budget balance removes the mean component from the planner's problem.

\begin{lemma}\label{lem:mean-invariance}
Suppose \(\omega_H=\omega_L\). Then the expected aggregate loss is invariant under the sharing rule, and the planner's welfare varies only through the variance terms.
\end{lemma}

\begin{proof}
Summing net losses across agents and using budget balance yields
\(\sum_{i=1}^n \tilde X_i = \sum_{i=1}^n X_i\).
Taking expectations gives invariance of the total expected loss. Under type-neutral welfare weights, this implies that the mean component of welfare is constant across feasible allocations.
\end{proof}

Under heterogeneous welfare weights, mean invariance need not hold at the planner level even though the sharing rule remains budget balanced in the aggregate.



\section{Independent Risks}\label{sec:ind}

We now specialize to the benchmark case in which losses are independent across agents. Under independence, the participation constraints become explicit, the feasible set admits a tractable geometry, and the planner's problem can be reduced to boundary optimization.

\subsection{Utility and feasibility under independence}

Assume that the losses of the agents are independent. Let
\[
T := n_H \alpha_H + n_L \alpha_L,
\qquad
P := n_H \alpha_H \mu_H + n_L \alpha_L \mu_L,
\qquad
V := n_H \alpha_H^2 \sigma_H^2 + n_L \alpha_L^2 \sigma_L^2 .
\]
Under truthful play, a representative type-\(k\) agent has mean net loss
\[
\mathbb E[\tilde X_k]
=
\mu_k-\alpha_k\!\left(\mu_k-\frac{P}{T}\right)
\]
and variance
\[
\Var(\tilde X_k)
=
\sigma_k^2(1-\alpha_k)^2 + 2\sigma_k^2(1-\alpha_k)\frac{\alpha_k^2}{T} + \alpha_k^2 \frac{V}{T^2}.
\]
Hence utility is
\begin{equation}\label{eq:utility-alpha-ind-new}
U_k(\alpha_H,\alpha_L)
=
-\mu_k
+
\alpha_k\!\left(\mu_k-\frac{P}{T}\right)
-
\frac{\gamma}{2}
\left[
\sigma_k^2(1-\alpha_k)^2 + 2\sigma_k^2(1-\alpha_k)\frac{\alpha_k^2}{T} + \alpha_k^2 \frac{V}{T^2}
\right].
\end{equation}



The autarkic allocation \((\alpha_H,\alpha_L)=(0,0)\) satisfies IR for both types with equality. We therefore restrict attention to profiles with \(T>0\).

For the high-risk type, the participation constraint is
\begin{equation}\label{eq:IRH-ind-new}
    \frac{\gamma_H}{2}\Big[
        2\sigma_H^2 - \alpha_H \sigma_H^2
        - 2(1-\alpha_H)\frac{\alpha_H}{T}\sigma_H^2
        - \frac{\alpha_H}{T^2} V
    \Big] \ge \frac{P}{T} - \mu_H.
\end{equation}
For the low-risk type, if \(\alpha_L>0\), the participation constraint is
\begin{equation}\label{eq:IRL-ind-new}
    \frac{\gamma_L}{2}\Big[
        2\sigma_L^2 - \alpha_L \sigma_L^2
        - 2(1-\alpha_L)\frac{\alpha_L}{T}\sigma_L^2
        - \frac{\alpha_L}{T^2} V
    \Big] \ge \frac{P}{T} - \mu_L.
\end{equation}
Clearing denominators yields
\begin{equation}\label{eq:IRLcore-new}
F_L(\alpha_H,\alpha_L)\ge 0,
\end{equation}
where
\begin{align}
F_L(\alpha_H,\alpha_L)
:=
\frac{\gamma_L}{2}\Big\{&
    (-n_L+1)n_L \sigma_L^2 \alpha_L^3
    + 2(-n_L+1)(n_H \alpha_H - n_L)\sigma_L^2 \alpha_L^2 \notag\\
    &+ \Big[(4n_L n_H - 2n_H - n_H^2 \alpha_H)\sigma_L^2 \alpha_H - \sigma_H^2 n_H \alpha_H^2\Big]\alpha_L
    + 2n_H^2 \sigma_L^2 \alpha_H^2 \Big\}
\notag\\
&
- n_H \alpha_H (\mu_H - \mu_L)(n_L \alpha_L + n_H \alpha_H).
\label{eq:FL-expanded-new}
\end{align}

Exact IC requires truthful type-contingent participation to dominate cross-type deviations:
\begin{subequations}\label{eq:IC-ind-new}
\begin{align}
U_H(\alpha_H\mid \alpha_H,\alpha_L)
&\ge
U_H(\alpha_L\mid \alpha_H,\alpha_L),
\\
U_L(\alpha_L\mid \alpha_H,\alpha_L)
&\ge
U_L(\alpha_H\mid \alpha_H,\alpha_L).
\end{align}
\end{subequations}
Closed-form expressions are reported in Appendix~\ref{app:IC}. For the main text, we maintain the sufficient conditions
\begin{equation}\label{eq:IC-conditions-new}
    \frac{\mu_H - \mu_L}{\gamma \sigma_L^2} \ge 
    \begin{cases}
        \tfrac{1}{2}, & \text{if } n_L < n_H,\\[6pt]
        \displaystyle \frac{(n_L + n_H)^2}{8 n_L n_H},
        & \text{if } n_L \ge n_H,
    \end{cases}
\end{equation}
which ensure truthful type-contingent participation.

\subsection{Geometry of the feasible set}

Define
\[
\mathcal S_{\mathrm{IR-L}}
:=
\bigl\{(\alpha_H,\alpha_L)\in[0,1]^2: F_L(\alpha_H,\alpha_L)\ge 0\bigr\}.
\]
The set of allocations satisfying the low-risk participation requirement is
\[
\mathcal S_{\mathrm{IR-L}}
\cup
\bigl\{(\alpha_H,0):\alpha_H\in[0,1]\bigr\}.
\]

\begin{proposition}\label{prop:IRL-properties-new}
The set \(\mathcal S_{\mathrm{IR-L}}\) contains the origin on its boundary and is
star-shaped with respect to the origin. Moreover, for each fixed
\(\alpha_L\in[0,1]\), the map \(\alpha_H\mapsto F_L(\alpha_H,\alpha_L)\) admits
at most one positive cutoff, so the low-risk participation constraint induces a
maximal feasible contribution \(\alpha_H^{\max}(\alpha_L)\in[0,1]\).
\end{proposition}

\begin{proof}
Since \(F_L(0,0)=0\), the origin belongs to the boundary of
\(\mathcal S_{\mathrm{IR-L}}=\{(\alpha_H,\alpha_L)\in[0,1]^2:F_L(\alpha_H,\alpha_L)\ge 0\}\).

To prove star-shapedness, decompose \(F_L=F_{L,2}+F_{L,3}\), where
\(F_{L,2}\) and \(F_{L,3}\) are the homogeneous components of total degree two
and three, respectively. For \((\alpha_H,\alpha_L)\in[0,1]^2\) with \(n_L\ge 1\),
one verifies that
\[
F_{L,3}(\alpha_H,\alpha_L)\le 0.
\]
Hence, for any \((\alpha_H,\alpha_L)\in \mathcal S_{\mathrm{IR-L}}\) and any
\(t\in[0,1]\),
\[
F_L(t\alpha_H,t\alpha_L)
= t^2 F_{L,2}(\alpha_H,\alpha_L)+t^3 F_{L,3}(\alpha_H,\alpha_L)
\ge t^2\bigl(F_{L,2}(\alpha_H,\alpha_L)+F_{L,3}(\alpha_H,\alpha_L)\bigr)
= t^2 F_L(\alpha_H,\alpha_L)\ge 0.
\]

Therefore \((t\alpha_H,t\alpha_L)\in\mathcal S_{\mathrm{IR-L}}\) for all
\(t\in[0,1]\), proving that \(\mathcal S_{\mathrm{IR-L}}\) is star-shaped with
respect to the origin.

Now fix \(\alpha_L\in[0,1]\). By \eqref{eq:FL-expanded-new},
\(F_L(\alpha_H,\alpha_L)\) is a quadratic polynomial in \(\alpha_H\). As shown
above, along the horizontal slice with this fixed \(\alpha_L\), the feasible set
cannot lose feasibility and then regain it at two distinct positive values of
\(\alpha_H\); equivalently, the equation
\[
F_L(\alpha_H,\alpha_L)=0
\]
admits at most one positive solution in \([0,1]\). It follows that the section
\[
\{\alpha_H\in[0,1]: F_L(\alpha_H,\alpha_L)\ge 0\}
\]
is of the form \([0,\alpha_H^{\max}(\alpha_L)]\) for some
\(\alpha_H^{\max}(\alpha_L)\in[0,1]\), which is the maximal feasible
high-risk contribution level induced by the low-risk participation constraint.
\end{proof}

\begin{remark}\label{rem:backward-bending}
Proposition~\ref{prop:IRL-properties-new} only implies the existence of a
slice-wise maximal feasible value \(\alpha_H^{\max}(\alpha_L)\). It does not
imply that the map \(\alpha_L\mapsto \alpha_H^{\max}(\alpha_L)\) is monotone.
In particular, the low-risk IR frontier may bend backward. Consequently, the
right boundary of the feasible set need not be monotone in the usual geometric
sense, even though each horizontal section is an interval.
\end{remark}

Define the feasible set
\[
\mathcal F
:=
\left\{(\alpha_H,\alpha_L) \in [0,1]^2 :
\text{\eqref{eq:IRH-ind-new}, \eqref{eq:IRL-ind-new}, and \eqref{eq:IC-ind-new} hold}
\right\}.
\]
Along any fixed direction \((\bar\alpha_H,\bar\alpha_L)\) with
\(\bar T=n_H\bar\alpha_H+n_L\bar\alpha_L>0\), write
\((\alpha_H,\alpha_L)=t(\bar\alpha_H,\bar\alpha_L)\). Substituting this ray
scaling into \eqref{eq:utility-alpha-ind-new} and collecting terms in \(t\)
gives
\[
W(t)=W(0)+A_1t-A_2t^2,
\]
where
\begin{align*}
A_1
={}&
\sum_{k\in\{H,L\}} n_k \omega_k \bar\alpha_k
\left(\mu_k-\frac{\bar P}{\bar T}\right)
+
\gamma
\sum_{k\in\{H,L\}} n_k \omega_k \bar\alpha_k \sigma_k^2,
\\
A_2
={}&
\frac{\gamma}{2}
\sum_{k\in\{H,L\}} n_k \omega_k
\left[
\bar\alpha_k^2 \sigma_k^2
+
\left(\frac{\bar\alpha_k}{\bar T}\right)^2 \bar V
\right].
\end{align*}


\begin{proposition}
\label{prop:boundary-opt-ind-new}
Suppose \(\omega_H=\omega_L\) and \(n_H,n_L\ge 2\). Then along every feasible
ray, welfare is weakly increasing on \([0,1]\). Consequently, every
welfare-maximizing feasible allocation lies on \(\partial\mathcal F\).
\end{proposition}

\begin{proof}
Fix any feasible ray and write welfare along the ray as
\[
W(t)=W(0)+\ell t-\frac{\gamma}{2}qt^2, \qquad t\in[0,1],
\]
for \(q\ge 0\), so \(W\) is concave in \(t\) as in Lemma~\ref{lem:concavity-general}.

Under type-neutral welfare weights, \(\omega_H=\omega_L\), and with
\(n_H,n_L\ge 2\), the independent benchmark implies that the ray-wise
maximizer weakly lies at or beyond \(t=1\). Equivalently,
\[
W'(1)=\ell-\gamma q \ge 0.
\]
Since \(W\) is concave, its derivative is weakly decreasing on \([0,1]\). Hence
for every \(t\in[0,1]\),
\[
W'(t)\ge W'(1)\ge 0.
\]
Therefore \(W\) is weakly increasing on \([0,1]\), so its maximum over the
feasible segment is attained at \(t=1\), i.e. at the boundary point of that ray.

Because every feasible allocation in \(\mathcal F\) lies on some feasible ray
emanating from the benchmark, it follows that any welfare-maximizing feasible
allocation must lie on \(\partial\mathcal F\).
\end{proof}


\subsection{Boundary characterization of the planner’s problem}

Define the right boundary
\[
\partial^+\mathcal F
:=
\left\{(\alpha_H,\alpha_L)\in\partial\mathcal F:
\nexists\,(\alpha_H',\alpha_L)\in\mathcal F
\text{ with }\alpha_H'>\alpha_H
\right\},
\]
and the upper feasible frontier
\[
\partial^\uparrow\mathcal F
:=
\left\{(\alpha_H,\alpha_L)\in\partial\mathcal F:
\nexists\,(\alpha_H,\alpha_L')\in\mathcal F
\text{ with }\alpha_L'>\alpha_L
\right\}.
\]

Under the conditions of Proposition~\ref{prop:boundary-opt-ind-new}, a welfare maximizer can be chosen
on the boundary of \(\mathcal F\). The next two propositions show that welfare
is weakly increasing along each vertical slice of the feasible set. Therefore,
for any given \(\alpha_H\), only the highest feasible \(\alpha_L\) is relevant
for welfare comparison, and the planner's problem can be reduced to
\(\partial^\uparrow\mathcal F\).

In the independent benchmark, \(\partial^+\mathcal F\) remains useful for
describing horizontal sections of the feasible region. However, the associated
boundary map \(\alpha_L\mapsto\alpha_H^{\max}(\alpha_L)\) need not be monotone,
so \(\partial^+\mathcal F\) may bend backward. Hence, absent an additional
monotonicity assumption, one cannot generally reduce the planner's problem to a
finite list of endpoint allocations.

\begin{proposition}\label{prop:vert-monotone-new}
Fix \(c\in(0,1]\) and consider the vertical slice
\[
\{(\alpha_H,\alpha_L): \alpha_H=c,\ 0\le \alpha_L\le c\}.
\]
Under type-neutral welfare weights, planner welfare is weakly increasing in
\(\alpha_L\) along this slice. Moreover, it is strictly increasing on \((0,c)\)
whenever \(n_L>1\) and \(\sigma_L^2>0\).
\end{proposition}

\begin{proof}
For fixed \(c\), define \(\widetilde W(\alpha_L):=W(c,\alpha_L)\) for
\(\alpha_L\in[0,c]\). Lemma~\ref{lem:vertical-type-neutral} in
Appendix~\ref{app:vertical} shows that
\[
\widetilde W'(\alpha_L)\ge 0
\qquad\text{for all }\alpha_L\in[0,c],
\]
with strict inequality for \(\alpha_L\in(0,c)\) whenever \(n_L>1\) and
\(\sigma_L^2>0\). This proves the claim.
\end{proof}

\begin{proposition}\label{prop:vert-monotone-general-omega}
Fix \(c\in(0,1]\) and consider the vertical slice
\[
\{(\alpha_H,\alpha_L): \alpha_H=c,\ 0\le \alpha_L\le c\}.
\]
If \(\omega_H\ge \omega_L\), then planner welfare is weakly increasing in
\(\alpha_L\) along this slice. Moreover, it is strictly increasing on \((0,c)\)
whenever \(n_L>1\) and \(\sigma_L^2>0\).
\end{proposition}

\begin{proof}
Write planner welfare as
\[
W(\alpha_H,\alpha_L)
=
\omega_L\bigl(n_H U_H(\alpha_H,\alpha_L)+n_L U_L(\alpha_H,\alpha_L)\bigr)
+(\omega_H-\omega_L)n_H U_H(\alpha_H,\alpha_L).
\]
Along the slice \(\alpha_H=c\), the first term is planner welfare under
type-neutral weights, which is weakly increasing in \(\alpha_L\) by
Proposition~\ref{prop:vert-monotone-new}. It therefore suffices to show that
\(U_H(c,\alpha_L)\) is weakly increasing in \(\alpha_L\).

Appendix~\ref{app:vertical} shows that
\[
\frac{\partial U_H(\alpha_H,\alpha_L)}{\partial \alpha_L}\ge 0
\]
throughout the feasible region. Hence, \(U_H(c,\alpha_L)\) is weakly increasing
in \(\alpha_L\), and so is planner welfare whenever \(\omega_H\ge\omega_L\).
The strict interior monotonicity follows from
Proposition~\ref{prop:vert-monotone-new} under \(n_L>1\) and \(\sigma_L^2>0\).
\end{proof}

\begin{proposition}
\label{prop:boundary-upper-general}
Suppose \(\omega_H=\omega_L\) and \(\mathcal F\neq\emptyset\). Then there exists
a welfare-maximizing allocation on the upper feasible frontier of
\(\mathcal F\). Equivalently, the planner's problem reduces to comparing
welfare across slice-wise highest feasible allocations.
\end{proposition}

\begin{proof}
By Proposition~\ref{prop:boundary-opt-ind-new}, there exists a
welfare-maximizing allocation on the boundary of \(\mathcal F\). Consider any
such maximizer \((\alpha_H,\alpha_L)\). If it is not slice-wise highest
feasible at the given \(\alpha_H\), then there exists
\((\alpha_H,\alpha_L')\in\mathcal F\) with \(\alpha_L'>\alpha_L\). By
Proposition~\ref{prop:vert-monotone-new}, welfare is weakly increasing in
\(\alpha_L\) along each vertical slice, so
\(W(\alpha_H,\alpha_L')\ge W(\alpha_H,\alpha_L)\). Therefore,
\((\alpha_H,\alpha_L)\) can be replaced, without lowering welfare, by a
feasible allocation with the same \(\alpha_H\) and maximal feasible
\(\alpha_L\).

Hence there exists a welfare-maximizing allocation on the upper feasible
frontier of \(\mathcal F\). Equivalently, the planner's problem reduces to
comparing welfare across slice-wise highest feasible allocations.
\end{proof}

Taken together, Propositions~\ref{prop:vert-monotone-new},
\ref{prop:vert-monotone-general-omega}, and
\ref{prop:boundary-upper-general} imply that, for each fixed \(\alpha_H\),
welfare is maximized at the largest feasible value of \(\alpha_L\). Thus the
planner's problem reduces to comparing welfare on the upper feasible frontier
\(\partial^\uparrow\mathcal F\).

The geometry of the feasible set implies that the planner's problem can be
organized around the welfare-relevant portion of the boundary. When the induced
right-boundary map \(\alpha_H^{\max}(\alpha_L)\) is monotone, this comparison
further collapses to a small set of endpoint candidates:
\begin{enumerate}[leftmargin=2em]
    \item the corner \((1,0)\);
    \item the full-pooling point \((1,1)\), whenever feasible;
    \item the terminal point at which the right boundary meets the low-risk
    participation frontier.
\end{enumerate}

Without such an
additional monotonicity property, however, the planner must in general compare
allocations along the upper-boundary $\partial^\uparrow\mathcal F$.

The feasibility of the full-pooling allocation \((1,1)\) is characterized by
the low-risk participation constraint. Appendix~\ref{app:pooling-feasibility}
derives the exact condition, provides a simple sufficient condition, and shows
that in large economies the exact threshold converges to
\[
    \frac{\mu_H - \mu_L}{\gamma \sigma_L^2}
    \le
    \frac{n_L + n_H}{2 n_H}.
\]




\subsection{Economic implications}

The previous discussion implies that, in large economies, the feasibility of
full pooling is governed by a simple asymptotic threshold. When population
shares converge to strictly interior limits, the finite-population correction
in the exact feasibility condition vanishes, so the condition for \((1,1)\)
approaches its large-pool counterpart.

\begin{proposition}
Suppose that \(n_H,n_L \to \infty\) and that the population shares converge to
strictly interior limits. Then the exact feasibility condition for the
full-pooling allocation \((1,1)\) converges to
\[
    \frac{\mu_H-\mu_L}{\gamma\sigma_L^2}
    \le
    \frac{n_L+n_H}{2n_H}.
\]
Equivalently, full pooling is feasible for all sufficiently large populations
whenever this limiting inequality holds strictly.
\end{proposition}




\section{Correlated Risks: How the Analysis Extends}
\label{sec:correlated}

We now discuss how the preceding analysis extends when losses are correlated
across agents. When losses are correlated, the independent-case characterization no longer goes through. We isolate the objects whose behavior must be controlled and the additional restrictions needed to recover analogous conclusions.

\subsection{Type-symmetric covariance structure and participation constraints}

To obtain tractable type-level formulas, suppose that covariance depends only on
types. For any distinct agents \(i\neq j\),
\begin{align*}
\Var(X_i)&=v_k
&&\text{if }\tau(i)=k,\\
\Cov(X_i,X_j)&=c_{kk}
\qquad&&\text{if }\tau(i)=\tau(j)=k,\\
\Cov(X_i,X_j)&=c_{HL}
\qquad&&\text{if }\tau(i)\neq\tau(j)&.
\end{align*}

Define the aggregate pooled loss
\[
Q:=\sum_{j=1}^n \alpha_{\tau(j)}X_j.
\]
Then the realized variance of a type-\(k\) agent's net loss can be written as
\[
\Var(\tilde X_k)
=
(1-\alpha_k)^2v_k
+
\left(\frac{\alpha_k}{T}\right)^2V_Q
+
2(1-\alpha_k)\frac{\alpha_k}{T}\Gamma_k,
\]
where
$V_Q:=\Var(Q)$ and $\Gamma_k:=\Cov(X_k,Q)$.

Under type symmetry,
\begin{align*}
V_Q
&=
\alpha_H^2\bigl[n_Hv_H+n_H(n_H-1)c_{HH}\bigr]
+
\alpha_L^2\bigl[n_Lv_L+n_L(n_L-1)c_{LL}\bigr]
\notag\\
&\quad
+
2\alpha_H\alpha_Ln_Hn_Lc_{HL},
\label{eq:VQ-corr-new}
\\
\Gamma_H
&=
\alpha_H\bigl[v_H+(n_H-1)c_{HH}\bigr]+\alpha_Ln_Lc_{HL},
\\
\Gamma_L
&=
\alpha_L\bigl[v_L+(n_L-1)c_{LL}\bigr]+\alpha_Hn_Hc_{HL}.
\end{align*}
Thus the correlated case differs from the independent benchmark through two
additional channels: covariance changes pooled risk \(V_Q\), and it also
changes each type's exposure to the pool through \(\Gamma_H\) and \(\Gamma_L\).

%\subsection{Participation constraints under covariance}

The high-risk and low-risk participation constraints become
\begin{equation}\label{eq:IRH-corr-new}
\alpha_H\!\left(\mu_H-\frac{P}{T}\right)
\ge
\frac{\gamma}{2}
\left[
(1-\alpha_H)^2v_H
+
\left(\frac{\alpha_H}{T}\right)^2V_Q
+
2(1-\alpha_H)\frac{\alpha_H}{T}\Gamma_H
-
v_H
\right],
\end{equation}
and
\begin{equation}\label{eq:IRL-corr-new}
\alpha_L\!\left(\mu_L-\frac{P}{T}\right)
\ge
\frac{\gamma}{2}
\left[
(1-\alpha_L)^2v_L
+
\left(\frac{\alpha_L}{T}\right)^2V_Q
+
2(1-\alpha_L)\frac{\alpha_L}{T}\Gamma_L
-
v_L
\right].
\end{equation}

\begin{remark}
Relative to the independent benchmark, feasibility now depends not only on
retained variance and pooled variance, but also on the covariance between an
agent's own loss and the aggregate pool. This covariance-exposure term is the
key new object generated by dependence.
\end{remark}

\subsection{What survives from the general theory}

The quadratic welfare representation (Proposition 3.1), the concavity of welfare along rays (Lemma 3.2), and the conditional boundary principle (Proposition 3.3) continue to apply. These results continue to provide the basic architecture of the problem.
In particular, once the feasible set is known to be star-shaped and welfare is
weakly increasing along feasible rays, the planner's optimum can again be
sought on an appropriate boundary.

Verifying the assumptions of that boundary argument is no longer immediate. In the correlated case,
\begin{itemize}[leftmargin=2em]
    \item star-shapedness of the feasible set is no longer automatic;
    \item the low-risk participation frontier need not admit the same simple
    cutoff structure as under independence;
    \item the sign of marginal welfare effects depends on covariance exposure
    as well as on variance reduction.
\end{itemize}

The derivative formulas underlying this decomposition are collected in
Appendix~\ref{app:alphaL}.

\subsection{Further analysis}

Under the type-symmetric covariance parameterization above, the correlated case
can be analyzed in three steps.

First, substitute \eqref{eq:VQ-corr-new} and the expression for \(\Gamma_L\)
into \eqref{eq:IRL-corr-new} to obtain a reduced-form description of the
low-risk participation constraint,
\[
F_L^{\mathrm{cov}}(\alpha_H,\alpha_L)\ge 0.
\]
This reduced form summarizes how feasibility depends jointly on pooled variance
and covariance exposure. Its explicit polynomial representation may be reported
either in the main text or in an appendix, depending on the desired level of
detail.

Second, study the geometry of the feasible set implied by
\(F_L^{\mathrm{cov}}(\alpha_H,\alpha_L)\ge 0\). The key objective is to
identify conditions under which the feasible region preserves enough structure
to support boundary-based arguments---for example, star-shapedness or a
monotone right boundary. Such conditions would provide an analogue of the
independent-case reduction to boundary comparisons.

Third, analyze welfare variation along the relevant frontier. Under type-neutral
welfare weights, mean invariance continues to simplify the planner's problem,
but the welfare derivative now reflects three distinct forces:
\begin{enumerate}[leftmargin=2em]
    \item direct retained-risk effects;
    \item pooled-variance effects through \(V_Q\);
    \item covariance-exposure effects through \(\Gamma_H\) and \(\Gamma_L\).
\end{enumerate}
This decomposition makes clear exactly where correlation modifies the
independent-case logic and what additional restrictions are needed to recover
analogous monotonicity or boundary-optimality conclusions.

%\subsection{Interpretation}

The correlated model shows why the independent benchmark cannot simply be
extrapolated to general dependence. When losses are correlated, broader
participation need not generate the same diversification gains, because
additional participants also bring covariance exposure into the pool.
Cross-type dependence \(c_{HL}\) is especially important: it directly affects
how low-risk participation changes each type's exposure to aggregate risk.

For this reason, the correlated case is best viewed as an extension in which
the general framework remains valid, but sharper results require additional
restrictions on the covariance structure or a numerical analysis of the
resulting frontier.

\section{Planner Allocation and Decentralized Participation}
\label{sec:decentralized}

This section contrasts the planner's allocation with the outcome of a
decentralized participation game in which each type chooses its own
contribution rate. Throughout, we restrict attention to symmetric strategies
within type: all type-\(H\) agents choose a common \(\alpha_H\in[0,1]\), and
all type-\(L\) agents choose a common \(\alpha_L\in[0,1]\). This restriction
is without loss because agents are ex ante identical within type.

Given \((\alpha_H,\alpha_L)\), the post-sharing loss of a type-\(k\) agent is
\[
\tilde X_k=(1-\alpha_k)X_k+\frac{\alpha_k}{T}Q,
\qquad
T=n_H\alpha_H+n_L\alpha_L,
\qquad
Q=\sum_{j=1}^n \alpha_{\tau(j)}X_j.
\]
The decentralized environment differs from the planner's problem only in the
assignment of decision rights over \((\alpha_H,\alpha_L)\).

\subsection{Decentralized equilibrium}

The objective of a representative type-\(k\) agent is
\[
U_k(\alpha_k,\alpha_{-k})
=
-\E[\tilde X_k]-\frac{\gamma}{2}\Var(\tilde X_k),
\qquad k\in\{H,L\},
\]
where, under the type-symmetric covariance parameterization of
Section~\ref{sec:correlated},
\[
\Var(\tilde X_k)
=
(1-\alpha_k)^2v_k
+
\left(\frac{\alpha_k}{T}\right)^{\!2}V_Q
+
2(1-\alpha_k)\frac{\alpha_k}{T}\Gamma_k.
\]

An interior best response of type \(k\) is characterized by the first-order
condition
\begin{equation}
\frac{\partial U_k}{\partial \alpha_k}
=
-\frac{\partial}{\partial \alpha_k}\E[\tilde X_k]
-\frac{\gamma}{2}\frac{\partial}{\partial \alpha_k}\Var(\tilde X_k)
=0.
\label{eq:decentralized-BR-FOC}
\end{equation}

A decentralized equilibrium is a pair
\((\alpha_H^{\mathrm{NE}},\alpha_L^{\mathrm{NE}})\in[0,1]^2\) such that
\[
\alpha_H^{\mathrm{NE}}\in\argmax_{\alpha_H\in[0,1]}
U_H(\alpha_H,\alpha_L^{\mathrm{NE}}),
\qquad
\alpha_L^{\mathrm{NE}}\in\argmax_{\alpha_L\in[0,1]}
U_L(\alpha_H^{\mathrm{NE}},\alpha_L).
\]

That is, equilibrium corresponds to the intersection of the two best-response correspondences. Each type internalizes only its own expected
transfer and its own retained risk, whereas the planner maximizes
\[
W(\alpha_H,\alpha_L)=n_HU_H(\alpha_H,\alpha_L)+n_LU_L(\alpha_H,\alpha_L).
\]
The two problems therefore generically yield distinct allocations.

\subsection{Optimality of full participation by the high-risk type}
\label{sec:decentralized-fullH}

We first characterize the conditions under which, given an arbitrary
low-risk participation rate \(\alpha_L=\ell\in[0,1]\), the high-risk type
optimally selects \(\alpha_H=1\). The analysis isolates two channels through
which a marginal increase in \(\alpha_H\) affects \(U_H\): the expected
transfer and the variance of the retained position.

\paragraph{Mean component.}
Differentiating the expected post-sharing loss yields
\[
\frac{\partial}{\partial \alpha_H}\E[\tilde X_H]
=
\frac{(n_H-T)(\mu_HT-P)}{T^2},
\qquad
T=n_H\alpha_H+n_L\ell.
\]
Since \(T\) is strictly increasing in \(\alpha_H\), each factor in the
numerator changes sign at most once, so \(\E[\tilde X_H]\) is unimodal in
\(\alpha_H\) and attains its minimum on \([0,1]\) at an endpoint. Evaluating
the derivative at \(\alpha_H=1\) gives
\begin{equation}
\left.\frac{\partial \E[\tilde X_H]}{\partial \alpha_H}\right|_{\alpha_H=1,\,\alpha_L=\ell}
=
-\,\frac{n_L^{\,2}\,\ell^{\,2}\,(\mu_H-\mu_L)}{(n_H+n_L\ell)^{2}}
\;\le\;0,
\label{eq:dEH-dalphaH-at1}
\end{equation}
with strict inequality whenever \(\ell>0\). Combined with the fact that
\(\E[\tilde X_H](1,\ell)<\E[\tilde X_H](0,\ell)\) under \(\mu_H>\mu_L\), it
follows that the mean component of \(U_H\) is globally minimized at
\(\alpha_H=1\).

\paragraph{Variance component.}
Define
\[
a_H:=v_H+(n_H-1)c_{HH},
\qquad
a_L:=v_L+(n_L-1)c_{LL}.
\]
Under the type-symmetric covariance parameterization,
\[
V_Q=\alpha_H^2 n_Ha_H+\alpha_L^2 n_La_L+2\alpha_H\alpha_L n_Hn_L c_{HL},
\qquad
\Gamma_H=\alpha_H a_H+\alpha_L n_L c_{HL}.
\]
Differentiating
\(\Var(\tilde X_H)=(1-\alpha_H)^2v_H+(\alpha_H/T)^2V_Q+2(1-\alpha_H)(\alpha_H/T)\Gamma_H\)
with respect to \(\alpha_H\) and evaluating at \(\alpha_H=1\),
\(\alpha_L=\ell\), the terms involving \((1-\alpha_H)\) vanish and one
obtains
\begin{equation}
\left.\frac{\partial \Var(\tilde X_H)}{\partial \alpha_H}\right|_{\alpha_H=1,\,\alpha_L=\ell}
=
\frac{2\ell\,n_L}{T^3}\bigl(V_Q-T\,\Gamma_H\bigr),
\qquad T=n_H+n_L\ell.
\label{eq:dVarH-dalphaH-at1-raw}
\end{equation}
A direct computation gives
\[
V_Q-T\,\Gamma_H \;=\; \ell\,n_L\Bigl[\ell\bigl(a_L-n_Lc_{HL}\bigr)
+n_Hc_{HL}-a_H\Bigr],
\]
so that \eqref{eq:dVarH-dalphaH-at1-raw} reduces to
\begin{equation}
\left.\frac{\partial \Var(\tilde X_H)}{\partial \alpha_H}\right|_{\alpha_H=1,\,\alpha_L=\ell}
=
\frac{2\,n_L^{\,2}\,\ell^{\,2}}{(n_H+n_L\ell)^{3}}\,
\Bigl(\ell\bigl(a_L-n_Lc_{HL}\bigr)+n_Hc_{HL}-a_H\Bigr).
\label{eq:dVarH-dalphaH-at1}
\end{equation}

\paragraph{Boundary first-order condition.}
Combining \eqref{eq:dEH-dalphaH-at1} and \eqref{eq:dVarH-dalphaH-at1},
\begin{equation}
\left.\frac{\partial U_H}{\partial \alpha_H}\right|_{\alpha_H=1,\,\alpha_L=\ell}
=
\frac{n_L^{\,2}\,\ell^{\,2}}{(n_H+n_L\ell)^{2}}
\left[
(\mu_H-\mu_L)
-
\frac{\gamma}{n_H+n_L\ell}\,
\Bigl(\ell(a_L-n_Lc_{HL})+n_Hc_{HL}-a_H\Bigr)
\right].
\label{eq:dUH-dalphaH-at1}
\end{equation}

\begin{proposition}\label{prop:alphaH1-bestresponse}
Suppose \(U_H(\,\cdot\,,\ell)\) is concave on \([0,1]\). Then \(\alpha_H=1\)
is a best response to \(\alpha_L=\ell\) if and only if
\begin{equation}\label{eq:star-condition}
(n_H+n_L\ell)\,(\mu_H-\mu_L)
\;\ge\;
\gamma\,\Bigl[\ell\,(a_L-n_Lc_{HL})+n_Hc_{HL}-a_H\Bigr].
\tag{$\star$}
\end{equation}
\end{proposition}

\begin{proof}
Concavity of \(U_H(\,\cdot\,,\ell)\) implies that \(\alpha_H=1\) is optimal on
\([0,1]\) if and only if the left derivative at \(\alpha_H=1\) is nonnegative.
By \eqref{eq:dUH-dalphaH-at1}, this is equivalent to the bracketed term being
nonnegative, which rearranges to \eqref{eq:star-condition}.
\end{proof}

\paragraph{A structural sufficient condition.}
The right-hand side of \eqref{eq:star-condition} is non-positive—so that
\eqref{eq:star-condition} holds for every \(\mu_H>\mu_L\)—whenever
\begin{equation}\label{eq:sufficient-structural}
\ell\,(a_L-n_L c_{HL})+n_H c_{HL}\;\le\;a_H.
\end{equation}
Under the model's natural parameter ordering, this inequality is implied, for
instance, by
\[
n_H\ge n_L,\qquad
c_{HH}\ge c_{LL},\qquad
c_{HL}\le \min\!\Bigl(\tfrac{a_H}{n_H},\,\tfrac{a_L}{n_L}\Bigr).
\]
The first two inequalities formalize the notion of a relatively concentrated
and internally correlated high-risk pool; the third bounds cross-type
covariance, so that low-risk participation cannot inject excessive covariance
exposure into the pool.

\subsection{Worst-case low-risk participation}\label{sec:worstcase-ell}

When the structural sufficient condition \eqref{eq:sufficient-structural}
fails, the question reduces to whether \eqref{eq:star-condition} holds at the
worst-case \(\ell\in[0,1]\). Normalizing the right-hand side of
\eqref{eq:star-condition} by \((n_H+n_L\ell)\), define
\[
g(\ell)
:=
\frac{\gamma\,(A\,\ell+B)}{n_H+n_L\ell},
\qquad
A:=a_L-n_L c_{HL},
\quad
B:=n_H c_{HL}-a_H.
\]
Then \eqref{eq:star-condition} is equivalent to
\((\mu_H-\mu_L)\ge g(\ell)\). Since
\[
g'(\ell)
=\frac{\gamma\,(A\,n_H-B\,n_L)}{(n_H+n_L\ell)^{2}}
=\frac{\gamma\bigl(n_H a_L+n_L a_H-2 n_H n_L c_{HL}\bigr)}
{(n_H+n_L\ell)^{2}},
\]
the sign of \(g'\) is constant on \([0,1]\), yielding the following
dichotomy.

\begin{corollary}\label{cor:worst-ell}
Let \(\Delta:=\mu_H-\mu_L\).
\begin{enumerate}[leftmargin=2em,label=(\roman*)]
\item If \(n_H a_L+n_L a_H\ge 2n_Hn_L c_{HL}\), then \(g\) is weakly
increasing on \([0,1]\) and the binding case is \(\ell=1\):
\[
\alpha_H=1\text{ is a best response for every }\ell\in[0,1]
\iff
\Delta\;\ge\;\frac{\gamma\bigl(a_L-n_L c_{HL}+n_Hc_{HL}-a_H\bigr)}{n_H+n_L}.
\]
\item If \(n_H a_L+n_L a_H< 2n_Hn_L c_{HL}\), then \(g\) is strictly
decreasing and the binding case is \(\ell=0\):
\[
\alpha_H=1\text{ is a best response for every }\ell\in[0,1]
\iff
\Delta\;\ge\;\frac{\gamma\,(n_Hc_{HL}-a_H)}{n_H}.
\]
\end{enumerate}
\end{corollary}

Case (i) covers all configurations with sufficiently small cross-type
covariance \(c_{HL}\); case (ii) arises only under unusually strong positive
cross-type dependence.

\subsection{Specialization to the independent benchmark}

Setting \(c_{HH}=c_{LL}=c_{HL}=0\) yields \(a_H=v_H\), \(a_L=v_L\),
\(A=v_L\), and \(B=-v_H\), so that \eqref{eq:star-condition} reduces to
\begin{equation}
(n_H+n_L\ell)\,(\mu_H-\mu_L)\;\ge\;\gamma\,(\ell\,v_L-v_H).
\label{eq:star-ind}
\end{equation}
Since \(v_H>v_L>0\), the right-hand side of \eqref{eq:star-ind} is bounded
above by \(\gamma(v_L-v_H)<0\), while the left-hand side is nonnegative.
Hence \eqref{eq:star-ind} holds identically.

\begin{corollary}\label{cor:ind-alphaH1}
In the independent benchmark, \(\alpha_H=1\) is a best response of the
high-risk type for every \(\alpha_L\in[0,1]\). Consequently, every
decentralized equilibrium satisfies \(\alpha_H^{\mathrm{NE}}=1\), and any
divergence between the planner and the decentralized allocation is confined to
the low-risk type.
\end{corollary}

This conclusion is consistent with the boundary geometry developed in
Section~\ref{sec:ind}: the high-risk type's expected transfer gain and its
diversification gain both favor maximal pooling, so the two effects reinforce
one another. The economic content of decentralized inefficiency in the
independent benchmark therefore lies entirely on the low-risk side: low-risk
agents fail to internalize the diversification benefit that their
participation confers on high-risk agents.

\subsection{Decomposition of the participation wedge}

The wedge between the decentralized and planner allocations stems from
cross-type spillovers in participation. A marginal increase in \(\alpha_k\)
alters not only type \(k\)'s retained risk, but also the size and covariance
structure of the common pool, and therefore the insurance terms faced by the
opposite type. Decentralized choice fails to internalize these spillovers.

Formally, an interior planner optimum in \(\alpha_L\) satisfies
\[
\frac{\partial W}{\partial \alpha_L}
=
n_H\frac{\partial U_H}{\partial \alpha_L}
+
n_L\frac{\partial U_L}{\partial \alpha_L}
=0,
\]
whereas an interior decentralized choice by type \(L\) satisfies only
\[
\frac{\partial U_L}{\partial \alpha_L}=0.
\]
The omitted term is the cross-type spillover
\[
\mathcal{E}_L
:=
n_H\,\frac{\partial U_H}{\partial \alpha_L},
\]
with an analogous expression \(\mathcal{E}_H:=n_L\,\partial U_L/\partial\alpha_H\)
for type \(H\). The sign and magnitude of \(\mathcal E_H\) and
\(\mathcal E_L\) determine whether decentralized participation is
inefficiently low or inefficiently high relative to the planner allocation.

Combining this decomposition with Corollary~\ref{cor:ind-alphaH1} sharpens
the conclusion in the independent benchmark. Since
\(\alpha_H^{\mathrm{NE}}=1\), the wedge reduces to the single spillover
\(\mathcal E_L\) evaluated at \((\alpha_H,\alpha_L)=(1,\alpha_L^{\mathrm{NE}})\).
By Lemma~\ref{lem:UH-vertical-monotone}, \(\partial U_H/\partial\alpha_L\ge 0\),
so \(\mathcal E_L\ge 0\), and decentralized participation by the low-risk
type is therefore weakly below the planner level. Under dependence,
\(\partial U_H/\partial\alpha_L\) need not have a definite sign, so the wedge
takes the form of a genuine composition externality rather than a
unidirectional bias.

\subsection{Characterization under type-symmetric covariance}

Substituting the type-symmetric expressions for \(V_Q\), \(\Gamma_H\), and
\(\Gamma_L\) into each type's first-order condition produces a pair of
best-response functions, \(\alpha_H=B_H(\alpha_L)\) and
\(\alpha_L=B_L(\alpha_H)\), which are nonlinear in general. Closed-form
solutions are available only in special cases such as the independent
benchmark of Corollary~\ref{cor:ind-alphaH1}. In the general type-symmetric
case, the decentralized equilibrium is obtained by solving
\eqref{eq:star-condition} jointly with its low-risk analogue, and is
compared with the planner allocation through the spillover terms
\(\mathcal E_H\) and \(\mathcal E_L\). The numerical exercises in
Section~\ref{sec:numerical} carry out this comparison and quantify the
resulting participation wedge.


\subsection{Characterization under type-symmetric covariance}

Substituting the type-symmetric expressions for \(V_Q\), \(\Gamma_H\), and
\(\Gamma_L\) into each type's first-order condition yields the best-response
functions \(\alpha_H=B_H(\alpha_L)\) and \(\alpha_L=B_L(\alpha_H)\). These are
nonlinear in general, so closed-form solutions are available only in special
cases such as the independent benchmark of
Corollary~\ref{cor:ind-alphaH1}. Otherwise, the decentralized equilibrium can
be computed numerically using \eqref{eq:star-condition} and its low-risk
analogue, and compared directly with the planner allocation. The numerical
illustrations in Section~\ref{sec:numerical} report this comparison and
quantify the resulting participation wedge.


\section{Numerical illustrations}\label{sec:numerical}

This section illustrates the main economic forces of the model through a sequence of numerical examples. We report a series of numerical exercises that illustrate the feasible set, the participation frontier, the planner's allocation, and the gap between decentralized and planner-selected participation.

Our benchmark calibration has two risk types, \(H\) and \(L\), with equal population sizes \(n_H=n_L=20\), risk aversion \(\gamma=2\), and welfare weights \(\omega_H=\omega_L=1\). Unless otherwise stated, the low-risk type has mean \(\mu_L=1\) and variance \(\sigma_L^2=1\), while the high-risk type has mean \(\mu_H=1.15\) and variance \(\sigma_H^2=4\). In each exercise, we evaluate utilities on a fine grid over \((\alpha_H,\alpha_L)\in[0,1]^2\), impose participation relative to autarky, and then characterize the feasible set and the welfare-maximizing allocation.

\subsection{Independent risks: feasible pooling and the participation frontier}

We begin with the benchmark case in which risks are independent across agents. This case isolates the basic trade-off between diversification gains and the redistribution implied by pooling types with different mean risks.

Figure \ref{fig:boundary-cases} highlights the allocations that matter most for the planner's comparison along the economically relevant boundary of the feasible set. In the independent benchmark, the planner optimum lies on, or very close to, the participation frontier rather than in the interior. Intuitively, the planner has a monotone incentive to enlarge the pool as long as additional participation remains acceptable to the marginal low-risk agent. Once that participation constraint binds, further expansion is no longer feasible.

This geometry implies that three boundary allocations are of primary interest:
\begin{enumerate}[leftmargin=2em]
    \item the corner \((1,0)\);
    \item the full-pooling allocation \((1,1)\), whenever it is feasible;
    \item the terminal point at which the right boundary \(\partial^+\mathcal F\) intersects the low-risk participation frontier.
\end{enumerate}
Taken together, these allocations summarize the relevant planner comparison in the independent benchmark.

\begin{figure}[H]
    \centering

    \begin{subfigure}[b]{0.32\textwidth}
        \centering
        \includegraphics[width=\textwidth]{1.png}
        \caption{Full pooling at \((1,1)\)}
        \label{fig:boundary-case-fullpool}
    \end{subfigure}
    \hfill
    \begin{subfigure}[b]{0.32\textwidth}
        \centering
        \includegraphics[width=\textwidth]{2.png}
        \caption{Terminal point on \(\partial^+\mathcal F\)}
        \label{fig:boundary-case-terminal}
    \end{subfigure}
    \hfill
    \begin{subfigure}[b]{0.32\textwidth}
        \centering
        \includegraphics[width=\textwidth]{3.png}
        \caption{Corner solution \((1,0)\)}
        \label{fig:boundary-case-corner}
    \end{subfigure}

    \caption{Boundary allocations of primary interest in the independent benchmark. Panel (a) shows the full-pooling allocation \((1,1)\), panel (b) the terminal point at which the right boundary \(\partial^+\mathcal F\) meets the low-risk participation frontier, and panel (c) the corner allocation \((1,0)\).}
    \label{fig:boundary-cases}
\end{figure}

A natural next question is when full pooling is feasible. To study this, we vary the mean gap by setting
\[
\mu_H=\mu_L+\Delta, \qquad \Delta\in[0,2],
\]
holding all other parameters fixed. For each \(\Delta\), we compute the two participation slacks at full pooling,
\[
U_L(1,1)-U_L^0
\qquad\text{and}\qquad
U_H(1,1)-U_H^0.
\]

\begin{figure}[H]
    \centering
    \includegraphics[width=0.72\textwidth]{figure_Y_full_pooling_slacks.pdf}
    \caption{Participation slacks at full pooling as the mean-risk gap \(\Delta\) increases. The figure plots \(U_L(1,1)-U_L^0\) and \(U_H(1,1)-U_H^0\) as functions of \(\Delta\), holding all other benchmark parameters fixed.}
    \label{fig:full-pooling-slacks}
\end{figure}

Figure \ref{fig:full-pooling-slacks} shows that full pooling is not automatically feasible. As \(\Delta\) increases, the low-risk participation slack declines monotonically in the benchmark calibration, while the high-risk slack remains positive over a substantially wider range. Hence the relevant threshold is determined by the low-risk participation constraint. Economically, feasibility at \((1,1)\) depends on whether the variance reduction from broader pooling is large enough to offset the expected redistribution toward the high-risk type. In this calibration, the low-risk type is the margin along which full pooling eventually breaks down.

To further illustrate the boundary nature of the planner solution, we also consider a more restrictive calibration with \(\mu_H=1.35\), keeping all other benchmark parameters unchanged. In that case, the low-risk participation frontier shifts inward and the planner optimum moves along the boundary rather than remaining near full participation. The planner's optimum thus tracks the low-risk participation frontier whenever the frontier binds.

\subsection{Dependence and covariance exposure}

We next introduce dependence across risks. The key difference from the independent benchmark is that participation now affects not only the size of the pool, but also its covariance structure. As a result, the marginal participant can either weaken or strengthen aggregate risk sharing depending on the pattern of covariance exposure.

To keep the calibration transparent while preserving positive semidefiniteness, we use the factor structure
\[
X_{ki}=\mu_k+b_k Z+d_k Z_k+\varepsilon_{ki},
\]
where \(Z\) is a common factor, \(Z_k\) is a type-specific factor, and \(\varepsilon_{ki}\) is idiosyncratic. This implies
\[
c_{HH}=b_H^2+d_H^2,\qquad
c_{LL}=b_L^2+d_L^2,\qquad
c_{HL}=b_H b_L.
\]
We choose idiosyncratic variances so that total variances remain fixed at \(v_H=4\) and \(v_L=1\), matching the independent benchmark.

Our first correlated calibration introduces positive common exposure:
\[
b_H=1.2,\qquad b_L=0.6,\qquad d_H=0.4,\qquad d_L=0.2,
\]
so that
\[
c_{HH}=1.60,\qquad c_{LL}=0.40,\qquad c_{HL}=0.72.
\]
For each \((\alpha_H,\alpha_L)\), we compute the relevant variance terms in the sharing arrangement, including \(V_Q\) and the exposure terms \(\Gamma_H,\Gamma_L\), and then evaluate utilities, participation, and welfare exactly as in the benchmark case.

\begin{figure}[H]
    \centering
    \includegraphics[width=0.95\textwidth]{figure_Z_correlated_vs_independent.pdf}
    \caption{Comparison of the independent benchmark and a positive-covariance calibration. Positive covariance reduces the diversification gain from additional participation and shifts the feasible set and planner optimum inward.}
    \label{fig:correlated-vs-independent}
\end{figure}

Relative to the independent benchmark, positive covariance reduces the diversification benefit of expanding the pool. Accordingly, the feasible set contracts, or rotates inward, in the direction relevant for additional low-risk participation, and the planner optimum also shifts inward. The intuition is straightforward: when a marginal participant brings risk that is positively exposed to the risks already in the pool, that participant contributes less diversification than under independence.

Our second correlated calibration introduces negative cross-type covariance, so that the two types partially hedge one another:
\[
b_H=1.2,\qquad b_L=-0.4,\qquad d_H=0.4,\qquad d_L=0.2,
\]
which implies
\[
c_{HH}=1.60,\qquad c_{LL}=0.20,\qquad c_{HL}=-0.48.
\]
Again, idiosyncratic variances are chosen so that total variances remain fixed at \(4\) and \(1\).

\begin{figure}[H]
    \centering
    \includegraphics[width=0.72\textwidth]{figure_ZZ_three_feasible_sets.pdf}
    \caption{Feasible sets under three covariance structures: the independent benchmark, a positive-covariance calibration, and a negative cross-type covariance calibration. Positive covariance shrinks feasible pooling, whereas negative cross-type covariance can enlarge it by making participation a source of hedging.}
    \label{fig:three-feasible-sets}
\end{figure}

Figure \ref{fig:three-feasible-sets} makes clear that correlation is not merely a friction. Positive common exposure tends to shrink feasible pooling because it weakens diversification, whereas negative cross-type covariance can enlarge the feasible set by making additional participation valuable as a hedge. Thus, once risks are dependent, the marginal value of participation becomes explicitly covariance-dependent.



\subsection{Planner and decentralized participation}

Our final exercise compares the planner allocation to decentralized participation. The objective is to visualize the wedge generated by cross-type spillovers.

We return to the independent benchmark and increase the mean gap slightly to
$\mu_H=1.20$,
so as to sharpen the distinction between planner and decentralized outcomes. For each \(\alpha_L\in[0,1]\), we compute the high-risk best response
\[
B_H(\alpha_L)\in \argmax_{\alpha_H\in[0,1]} U_H(\alpha_H,\alpha_L),
\]
and for each \(\alpha_H\in[0,1]\), the low-risk best response
\[
B_L(\alpha_H)\in \argmax_{\alpha_L\in[0,1]} U_L(\alpha_H,\alpha_L).
\]
Decentralized equilibria are then the intersections of the two best-response correspondences.

\begin{figure}[H]
    \centering
    \includegraphics[width=0.72\textwidth]{figure_W_planner_vs_decentralized.pdf}
    \caption{Planner allocation versus decentralized participation in the independent benchmark with \(\mu_H=1.20\). The figure plots the two best-response curves together with the feasible region and the planner-selected allocation.}
    \label{fig:planner-vs-decentralized}
\end{figure}

Figure \ref{fig:planner-vs-decentralized} shows that decentralized participation generally differs from the planner optimum. The decentralized equilibrium satisfies each type's private marginal condition, but not the planner's marginal condition, because each type ignores the effect of its own participation on the other type's risk-sharing opportunities. In the benchmark calibration, the low-risk type participates less in decentralized equilibrium than under the planner allocation. The reason is that low-risk agents do not internalize the diversification benefit their participation confers on high-risk agents. The planner does internalize this spillover and therefore selects a point with [TO BE FILLED: typically higher \(\alpha_L\), possibly also different \(\alpha_H\)].

To show that even the sign of the participation wedge can depend on covariance structure, we repeat the comparison under the two correlated calibrations introduced above. Table X reports the planner allocation, the decentralized equilibrium, the welfare gap, and the spillover term
\[
n_{-k}\frac{\partial U_{-k}}{\partial \alpha_k}
\]
evaluated at equilibrium for each type \(k\).

[Table X about here]

The table makes clear that the wedge is not mechanically an underparticipation result for one particular type. Under positive common exposure, the planner may prefer to restrain participation relative to the independent benchmark because marginal entrants contribute less diversification and more covariance exposure. Under negative cross-type covariance, by contrast, the planner may encourage broader participation because one type's entry improves the hedging opportunities of the other. Accordingly, the decentralized distortion is best understood not as a fixed type-based bias, but as a composition externality: one type's participation affects the other's post-sharing variance, and decentralized choice fails to internalize that effect.

The numerical exercises show that the participation frontier shapes feasible allocations even under independence; that under dependence the effect of additional participation hinges on covariance exposure, with positive common exposure shrinking and negative cross-type exposure enlarging the feasible set; and that decentralized participation departs from the planner allocation through a composition externality whose sign depends on the covariance structure.


\section{Conclusion}\label{sec:conclusion}

This paper studies participation and optimal design in peer-to-peer
insurance with heterogeneous agents. The central point is that
heterogeneity matters not only because agents face different risks,
but also because different risk categories may make different
participation decisions. The viability of a P2P arrangement therefore
depends on the endogenous composition of the pool, rather than on its
aggregate scale alone.

We develop a general framework in which heterogeneous agents choose
type-contingent participation rates under a common budget-balanced
sharing rule. The planner's problem admits a tractable geometric
structure: welfare can be represented quadratically along rays, and
optimal allocations satisfy a conditional boundary principle. These
results identify the feasibility and welfare objects that determine
when pooling arrangements can improve upon outside insurance options.

In the independent-risk benchmark, this structure yields sharper
characterizations. The low-risk participation frontier is tractable,
the planner's problem reduces to boundary optimization, and the
feasibility of full pooling can be characterized by an explicit
threshold. This threshold becomes asymptotically tight in large
populations, while also highlighting the importance of finite-pool
participation constraints: convergence to pure premiums does not, by
itself, eliminate the participation problem.

Allowing for correlated risks introduces covariance exposure as an
additional determinant of both feasibility and welfare. Depending on
the sign and magnitude of cross-type dependence, correlation may
either contract or expand the feasible set. The general framework
continues to apply, although sharper conclusions require additional
restrictions on the covariance structure.

The main economic implication arises from the comparison between
planner and decentralized participation. For a broad class of
parameter configurations characterized in the paper, the decentralized
game has an equilibrium of the form
\((\alpha_H^{\mathrm{NE}},\alpha_L^{\mathrm{NE}})
=(1,\ell^{\mathrm{NE}})\), with \(\ell^{\mathrm{NE}}<1\). High-risk
agents participate fully, whereas low-risk agents underparticipate
because they do not internalize the diversification benefit that their
participation confers on the high-risk pool. The planner internalizes
this cross-type spillover and can move the allocation along the upper
feasible frontier toward full pooling. In particular, under the
maintained assumptions of the model, whenever full pooling
\((1,1)\) is feasible, the planner implements it and strictly
Pareto-improves upon the decentralized outcome.

Thus, in a non-trivial region of the parameter space, centralized
design transforms a partially participating equilibrium into a fully
participating pool. This transition from
\((1,\ell^{\mathrm{NE}})\) to \((1,1)\) is the main takeaway of the
paper. It identifies a concrete margin on which centralized design
outperforms decentralized participation and isolates the low-risk
participation constraint as the key object that practical P2P
insurance designs must address.

More broadly, the paper suggests that P2P insurance should be
understood through a participation-based lens. The relevant design
question is not only how losses are shared within a given pool, but
which risk categories are willing to enter the pool in the first
place, and which institutional arrangements can sustain the pool
composition required to realize diversification gains. The wedge
between \((1,\ell^{\mathrm{NE}})\) and \((1,1)\) captures the central
design margin at stake.





\appendix

\section{Detailed Incentive Compatibility Derivations}\label{app:IC}

This appendix records the explicit algebraic expressions underlying the
incentive compatibility (IC) constraints used in the main text. All expressions
reported below are algebraically identical to those in the original working
draft; only explanatory text and signposting have been added. Simplifications
and factorizations were verified using symbolic algebra software to ensure
correctness. No new assumptions or results are introduced in this appendix.

\subsection{High-risk deviation to $\alpha_L$}\label{app:IC-H}

Let $T = n_H \alpha_H + n_L \alpha_L$ and define
\[
    T_H' := \alpha_H(n_H-1) + \alpha_L(n_L+1).
\]
The difference between remaining at $\alpha_H$ and deviating to $\alpha_L$ is
\begin{align}
    \Psi_H(\alpha_H,\alpha_L)
    &= U_H(\alpha_H \mid \alpha_H,\alpha_L)
       - U_H(\alpha_L \mid \alpha_H,\alpha_L).
\end{align}
After simplification,
\begin{align}
    \Psi_H(\alpha_H,\alpha_L)
    &= -\frac{1}{2}
       (\alpha_H n_H - \alpha_H + \alpha_L n_L)
       (\alpha_H - \alpha_L)\,
       \Xi_H(\alpha_H,\alpha_L), \label{eq:PsiH-app}
\end{align}
where the polynomial $\Xi_H(\alpha_H,\alpha_L)$ is given explicitly below.

\begin{align}
    \Xi_H(\alpha_H,\alpha_L)
    &= 2 \alpha_L^3 \mu_L n_L^2 - 2 \alpha_L^3 \mu_H n_L^3 - 2 \alpha_L^3 \mu_H n_L^2 + 2 \alpha_L^3 \mu_L n_L^3 \notag\\
    &\quad + 4 \alpha_H^3 \gamma n_H^2 \sigma_H^2 - 2 \alpha_H^3 \gamma n_H^3 \sigma_H^2 - 2 \alpha_H^4 \gamma n_H^2 \sigma_H^2 + \alpha_H^4 \gamma n_H^3 \sigma_H^2 \notag\\
    &\quad - 2 \alpha_L^3 \gamma n_L^2 \sigma_H^2 - 2 \alpha_L^3 \gamma n_L^3 \sigma_H^2 + \alpha_L^4 \gamma n_L^3 \sigma_H^2 + \alpha_L^4 \gamma n_L^2 \sigma_L^2 \notag\\
    &\quad + 2 \alpha_H \alpha_L^2 \mu_H n_L^2 - 2 \alpha_H \alpha_L^2 \mu_L n_L^2 - 2 \alpha_H^3 \gamma n_H \sigma_H^2 + \alpha_H^4 \gamma n_H \sigma_H^2 \notag\\
    &\quad + 2 \alpha_H^2 \alpha_L \gamma n_H \sigma_H^2 - \alpha_H^3 \alpha_L \gamma n_H \sigma_H^2 + 2 \alpha_H \alpha_L^2 \gamma n_L \sigma_H^2 \notag\\
    &\quad - 2 \alpha_H^2 \alpha_L \gamma n_L \sigma_H^2 + 2 \alpha_H^3 \alpha_L \gamma n_L \sigma_H^2 + \alpha_H \alpha_L^3 \gamma n_L \sigma_L^2 \notag\\
    &\quad - 4 \alpha_H \alpha_L^2 \mu_H n_H n_L^2 - 2 \alpha_H^2 \alpha_L \mu_H n_H^2 n_L + 4 \alpha_H \alpha_L^2 \mu_L n_H n_L^2 \notag\\
    &\quad + 2 \alpha_H^2 \alpha_L \mu_L n_H^2 n_L - 2 \alpha_H^2 \alpha_L \gamma n_H^2 \sigma_H^2 + \alpha_H^3 \alpha_L \gamma n_H^3 \sigma_H^2 \notag\\
    &\quad + 4 \alpha_H \alpha_L^2 \gamma n_L^2 \sigma_H^2 - 2 \alpha_H^2 \alpha_L^2 \gamma n_L \sigma_H^2 - \alpha_H \alpha_L^3 \gamma n_L^2 \sigma_H^2 \notag\\
    &\quad + \alpha_H \alpha_L^3 \gamma n_L^3 \sigma_H^2 - \alpha_H^2 \alpha_L^2 \gamma n_L \sigma_L^2 + \alpha_H \alpha_L^3 \gamma n_L^2 \sigma_L^2 \notag\\
    &\quad - 2 \alpha_H \alpha_L^2 \mu_H n_H n_L + 2 \alpha_H^2 \alpha_L \mu_H n_H n_L + 2 \alpha_H \alpha_L^2 \mu_L n_H n_L \notag\\
    &\quad - 2 \alpha_H^2 \alpha_L \mu_L n_H n_L - 3 \alpha_H^2 \alpha_L^2 \gamma n_L^2 \sigma_H^2 + 3 \alpha_H^2 \alpha_L^2 \gamma n_H n_L^2 \sigma_H^2 \notag\\
    &\quad + 3 \alpha_H^2 \alpha_L^2 \gamma n_H^2 n_L \sigma_H^2 - 4 \alpha_H \alpha_L^2 \gamma n_H n_L \sigma_H^2 + 8 \alpha_H^2 \alpha_L \gamma n_H n_L \sigma_H^2 \notag\\
    &\quad - 5 \alpha_H^3 \alpha_L \gamma n_H n_L \sigma_H^2 + \alpha_H \alpha_L^3 \gamma n_H n_L \sigma_L^2 \notag\\
    &\quad - 6 \alpha_H \alpha_L^2 \gamma n_H n_L^2 \sigma_H^2 - 6 \alpha_H^2 \alpha_L \gamma n_H^2 n_L \sigma_H^2 - \alpha_H^2 \alpha_L^2 \gamma n_H n_L \sigma_H^2 \notag\\
    &\quad + 3 \alpha_H \alpha_L^3 \gamma n_H n_L^2 \sigma_H^2 + 3 \alpha_H^3 \alpha_L \gamma n_H^2 n_L \sigma_H^2 + \alpha_H^2 \alpha_L^2 \gamma n_H n_L \sigma_L^2. \label{eq:XiH}
\end{align}

\paragraph{Special case $\alpha_H=1$.}
Setting $\alpha_H = 1$ in \eqref{eq:XiH} yields
\begin{align}
    \Xi_H(1,\alpha_L)
    &= (C_5 + C_4)\,\alpha_L^4 \gamma - (2 n_L^3 + 2 n_L^2)\,\alpha_L^3 \mu_H \notag\\
    &\quad + (2 n_L^3 + 2 n_L^2)\,\alpha_L^3 \mu_L + (C_4 - C_5 - 3 n_L^2 \sigma_H^2 + n_L \sigma_L^2 + C_7 + C_2)\,\alpha_L^3 \gamma \notag\\
    &\quad + (2 n_L^2 - C_8 - 2 n_H n_L)\,\alpha_L^2 \mu_H + (2 n_H n_L + C_8 - 2 n_L^2)\,\alpha_L^2 \mu_L \notag\\
    &\quad + (C_1 - C_2 - 5 n_H n_L \sigma_H^2 + C_7 + n_L^2 \sigma_H^2 - n_L \sigma_L^2)\,\alpha_L^2 \gamma \notag\\
    &\quad + (- C_9 + 2 n_L n_H)\,\alpha_L \mu_H + (C_9 - 2 n_L n_H)\,\alpha_L \mu_L \notag\\
    &\quad + (C_6 - C_3 + n_H \sigma_H^2 + 3 n_H n_L \sigma_H^2 - C_1)\,\alpha_L \gamma \notag\\
    &\quad + (- C_6 + C_3 - n_H \sigma_H^2)\,\gamma, \label{eq:XiH-special}
\end{align}

with
\begin{align*}
    C_1 &= 3 n_H^2 n_L \sigma_H^2, &
    C_2 &= 3 n_H n_L^2 \sigma_H^2, &
    C_3 &= 2 n_H^2 \sigma_H^2, \\
    C_4 &= n_L^2 \sigma_L^2, &
    C_5 &= n_L^3 \sigma_H^2, &
    C_6 &= n_H^3 \sigma_H^2, \\
    C_7 &= n_L n_H \sigma_L^2, &
    C_8 &= 4 n_H n_L^2, &
    C_9 &= 2 n_L n_H^2.
\end{align*}

\paragraph{Leading-order condition.}
Collecting the $n^3$ terms of \eqref{eq:XiH-special} gives
\begin{equation}\label{eq:IC-H-leading}
    (\mu_H - \mu_L)\Big[-2 n_L^3 \alpha_L^3 - 4 n_H n_L^2 \alpha_L^2 - 2 n_L n_H^2 \alpha_L \Big]
    + (\alpha_L - 1)(n_L \alpha_L + n_H)^3 \gamma \sigma_H^2 \le 0 .
\end{equation}

\subsection{Low-risk deviation to $\alpha_H$}\label{app:IC-L}

Let
\[
    T_L' := \alpha_H(n_H+1) + \alpha_L(n_L-1).
\]
Define
\begin{equation}
    \Psi_L(\alpha_H,\alpha_L)
    := U_L(\alpha_L \mid \alpha_H,\alpha_L)
       - U_L(\alpha_H \mid \alpha_H,\alpha_L).
\end{equation}
After simplification,
\begin{align}
    \Psi_L(\alpha_H,\alpha_L)
    = -\frac{1}{2}
      (\alpha_L n_L - \alpha_L + \alpha_H n_H)
      (\alpha_L - \alpha_H)\,
      \Xi_L(\alpha_H,\alpha_L), \label{eq:PsiL-app}
\end{align}

\begin{align}
    \Xi_L(\alpha_H,\alpha_L)
    &= 2 \alpha_H^3 \mu_H n_H^2 - 2 \alpha_H^3 \mu_L n_H^3 - 2 \alpha_H^3 \mu_L n_H^2 + 2 \alpha_H^3 \mu_H n_H^3 \notag\\
    &\quad + 4 \alpha_L^3 \gamma n_L^2 \sigma_L^2 - 2 \alpha_L^3 \gamma n_L^3 \sigma_L^2 - 2 \alpha_L^4 \gamma n_L^2 \sigma_L^2 + \alpha_L^4 \gamma n_L^3 \sigma_L^2 \notag\\
    &\quad - 2 \alpha_H^3 \gamma n_H^2 \sigma_L^2 - 2 \alpha_H^3 \gamma n_H^3 \sigma_L^2 + \alpha_H^4 \gamma n_H^3 \sigma_L^2 + \alpha_H^4 \gamma n_H^2 \sigma_H^2 \notag\\
    &\quad + 2 \alpha_L \alpha_H^2 \mu_L n_H^2 - 2 \alpha_L \alpha_H^2 \mu_H n_H^2 - 2 \alpha_L^3 \gamma n_L \sigma_L^2 \notag\\
    &\quad + \alpha_L^4 \gamma n_L \sigma_L^2 + 2 \alpha_L^2 \alpha_H \gamma n_L \sigma_L^2 - \alpha_L^3 \alpha_H \gamma n_L \sigma_L^2 \notag\\
    &\quad + 2 \alpha_L \alpha_H^2 \gamma n_H \sigma_L^2 - 2 \alpha_L^2 \alpha_H \gamma n_H \sigma_L^2 + 2 \alpha_L^3 \alpha_H \gamma n_H \sigma_L^2 \notag\\
    &\quad + \alpha_L \alpha_H^3 \gamma n_H \sigma_H^2 - 4 \alpha_L \alpha_H^2 \mu_L n_L n_H^2 - 2 \alpha_L^2 \alpha_H \mu_L n_L^2 n_H \notag\\
    &\quad + 4 \alpha_L \alpha_H^2 \mu_H n_L n_H^2 + 2 \alpha_L^2 \alpha_H \mu_H n_L^2 n_H - 2 \alpha_L^2 \alpha_H \gamma n_L^2 \sigma_L^2 \notag\\
    &\quad + \alpha_L^3 \alpha_H \gamma n_L^3 \sigma_L^2 + 4 \alpha_L \alpha_H^2 \gamma n_H^2 \sigma_L^2 - 2 \alpha_L^2 \alpha_H^2 \gamma n_H \sigma_L^2 \notag\\
    &\quad - \alpha_L \alpha_H^3 \gamma n_H^2 \sigma_L^2 + \alpha_L \alpha_H^3 \gamma n_H^3 \sigma_L^2 - \alpha_L^2 \alpha_H^2 \gamma n_H \sigma_H^2 \notag\\
    &\quad + \alpha_L \alpha_H^3 \gamma n_H^2 \sigma_H^2 - 2 \alpha_L \alpha_H^2 \mu_L n_L n_H + 2 \alpha_L^2 \alpha_H \mu_L n_L n_H \notag\\
    &\quad + 2 \alpha_L \alpha_H^2 \mu_H n_L n_H - 2 \alpha_L^2 \alpha_H \mu_H n_L n_H - 3 \alpha_L^2 \alpha_H^2 \gamma n_H^2 \sigma_L^2 \notag\\
    &\quad + 3 \alpha_L^2 \alpha_H^2 \gamma n_L n_H^2 \sigma_L^2 + 3 \alpha_L^2 \alpha_H^2 \gamma n_L^2 n_H \sigma_L^2 - 4 \alpha_L \alpha_H^2 \gamma n_L n_H \sigma_L^2 \notag\\
    &\quad + 8 \alpha_L^2 \alpha_H \gamma n_L n_H \sigma_L^2 - 5 \alpha_L^3 \alpha_H \gamma n_L n_H \sigma_L^2 + \alpha_L \alpha_H^3 \gamma n_L n_H \sigma_H^2 \notag\\
    &\quad - 6 \alpha_L \alpha_H^2 \gamma n_L n_H^2 \sigma_L^2 - 6 \alpha_L^2 \alpha_H \gamma n_L^2 n_H \sigma_L^2 - \alpha_L^2 \alpha_H^2 \gamma n_L n_H \sigma_L^2 \notag\\
    &\quad + 3 \alpha_L \alpha_H^3 \gamma n_L n_H^2 \sigma_L^2 + 3 \alpha_L^3 \alpha_H \gamma n_L^2 n_H \sigma_L^2 + \alpha_L^2 \alpha_H^2 \gamma n_L n_H \sigma_H^2. \label{eq:XiL}
\end{align}

\paragraph{Special case $\alpha_H=1$.}
Setting $\alpha_H = 1$ yields
\begin{align}
    \Xi_L(1,\alpha_L)
    &= (D_5 - D_3 + n_L \sigma_L^2)\,\alpha_L^4 \gamma \notag\\
    &\quad + (4 n_L^2 \sigma_L^2 - D_5 + D_9 - 3 n_L \sigma_L^2 - 5 n_H n_L \sigma_L^2 + D_2)\,\alpha_L^3 \gamma \notag\\
    &\quad + (D_{11} - 2 n_H n_L)\,\alpha_L^2 \mu_H + (-D_{11} + 2 n_H n_L)\,\alpha_L^2 \mu_L \notag\\
    &\quad + (D_1 - D_4 - D_2 + D_8 + 7 n_H n_L \sigma_L^2 - n_H \sigma_H^2 - 4 n_H \sigma_L^2 - D_3 + 2 n_L \sigma_L^2)\,\alpha_L^2 \gamma \notag\\
    &\quad + (2 n_H n_L + D_{10} - 2 n_H^2)\,\alpha_L \mu_H + (2 n_H^2 - D_{10} - 2 n_H n_L)\,\alpha_L \mu_L \notag\\
    &\quad + (D_7 + D_4 + D_6 + n_H \sigma_H^2 + D_9 + D_8 - 4 n_H n_L \sigma_L^2 - D_1)\,\alpha_L \gamma \notag\\
    &\quad + (2 n_H^3 + 2 n_H^2)\,\mu_H + (-2 n_H^3 - 2 n_H^2)\,\mu_L + (-D_6 + D_7 - 2 n_H^2 \sigma_L^2)\,\gamma, \label{eq:XiL-special}
\end{align}

where
\begin{align*}
    D_1 &= 3 n_H^2 n_L \sigma_L^2, &
    D_2 &= 3 n_H n_L^2 \sigma_L^2, &
    D_3 &= 2 n_L^2 \sigma_L^2, \\
    D_4 &= 3 n_H^2 \sigma_L^2, &
    D_5 &= n_L^3 \sigma_L^2, &
    D_6 &= n_H^3 \sigma_L^2, \\
    D_7 &= n_H^2 \sigma_H^2, &
    D_8 &= n_H n_L \sigma_H^2, &
    D_9 &= 2 n_H \sigma_L^2, \\
    D_{10} &= 4 n_H^2 n_L, &
    D_{11} &= 2 n_H n_L^2.
\end{align*}

\paragraph{Leading-order condition.}
The $n^3$ components give
\begin{equation}\label{eq:IC-L-leading}
    2 n_H (n_L \alpha_L + n_H)^2 (\mu_H - \mu_L)
    + (\alpha_L - 1)(n_L \alpha_L + n_H)^3 \gamma \sigma_L^2 \ge 0.
\end{equation}

\paragraph{General $\alpha_H \neq 1$.}
For general $(\alpha_H,\alpha_L)\in[0,1]^2$, the expressions
\eqref{eq:PsiH-app}--\eqref{eq:XiL} characterize the exact finite-$n$ IC region.
The main text relies only on the factorizations and sufficient conditions
derived from these expressions.



\section{Vertical monotonicity on slices}
\label{app:vertical}

This appendix proves Propositions~\ref{prop:vert-monotone-new} and
\ref{prop:vert-monotone-general-omega}.



\begin{lemma}
\label{lem:vertical-type-neutral}
Fix \(c\in(0,1]\) and define
\[
\widetilde W(\alpha_L):=W(c,\alpha_L), \qquad \alpha_L\in[0,c].
\]
Under type-neutral welfare weights,
\[
\widetilde W'(\alpha_L)\ge 0
\qquad\text{for all }\alpha_L\in[0,c].
\]
If in addition \(n_L>1\) and \(\sigma_L^2>0\), then
\[
\widetilde W'(\alpha_L)>0
\qquad\text{for all }\alpha_L\in(0,c).
\]
\end{lemma}

\begin{proof}
Fix \(c\in(0,1]\). Differentiating welfare along the slice \(\alpha_H=c\) gives
\[
\frac{d\widetilde W}{d\alpha_L}
= \omega \gamma
   \frac{\Omega_0(c) + \Omega_1(c)\alpha_L + \Omega_2(c)\alpha_L^2
        + \Omega_3(c)\alpha_L^3 + \Omega_4(c)\alpha_L^4}
        {(n_H c + n_L \alpha_L)^3},
\]
where
\[
\begin{aligned}
\Omega_0(c) &= n_L n_H^2 c^3\big(\sigma_H^2 + n_H \sigma_L^2\big),\\
\Omega_1(c) &= - n_H n_L c^2 \sigma_H^2\big[c(n_H + n_L) - n_L\big]
              - n_H^2 n_L c^2 \sigma_L^2\big[(1+n_H)c + 2 - 3 n_L\big],\\
\Omega_2(c) &= 3 n_H n_L \sigma_L^2\, c\, (n_L - 1)(n_L - n_H c),\\
\Omega_3(c) &= - n_L^2 \sigma_L^2\, (n_L - 1)(3 n_H c - n_L),\\
\Omega_4(c) &= - n_L^3 \sigma_L^2\, (n_L - 1).
\end{aligned}
\]

Since \(\omega>0\), \(\gamma>0\), and \(n_H c+n_L\alpha_L>0\), the sign of
\(\widetilde W'(\alpha_L)\) is the sign of the numerator. We decompose that
numerator as
\[
\sigma_H^2 Q_H(\alpha_L)+\sigma_L^2 Q_L(\alpha_L),
\]
where
\[
Q_H(\alpha_L)
=
n_L n_H^2 c^3
-
n_H n_L c^2\bigl[c(n_H+n_L)-n_L\bigr]\alpha_L,
\]
and
\[
Q_L(\alpha_L)
=
n_L\!\left[
(1-\alpha_L)(t c+\alpha_L)^3(n_L-1)
+t^2 c^2\bigl((1-c-t c)\alpha_L+t c\bigr)
\right],
\qquad t:=\frac{n_H}{n_L}.
\]

We first show \(Q_H(\alpha_L)\ge 0\). Since \(Q_H\) is affine in \(\alpha_L\),
it suffices to evaluate it at the upper endpoint \(\alpha_L=c\):
\[
Q_H(c)
=
n_L n_H^2 c^3
-
n_H n_L c^3\bigl[c(n_H+n_L)-n_L\bigr]
=
n_H n_L c^3(1-c)(n_H+n_L)
\ge 0.
\]
Because the slope of \(Q_H\) is weakly negative, it follows that
\(Q_H(\alpha_L)\ge Q_H(c)\ge 0\) for every \(\alpha_L\in[0,c]\).

Next consider \(Q_L(\alpha_L)\). For \(0\le \alpha_L\le c\le 1\), we have
\(1-\alpha_L\ge 0\), \(t c+\alpha_L\ge 0\), and \(n_L-1\ge 0\). Hence the first
term inside the brackets,
\[
(1-\alpha_L)(t c+\alpha_L)^3(n_L-1),
\]
is weakly nonnegative. For the second term, note that
\[
(1-c-t c)\alpha_L+t c
\ge
(1-c-t c)c+t c
=
c(1-c)(1+t)
\ge 0,
\]
where the inequality uses \(\alpha_L\le c\), and the final expression is
nonnegative because \(c\in(0,1]\) and \(t>0\). Therefore
\[
t^2 c^2\bigl((1-c-t c)\alpha_L+t c\bigr)\ge 0.
\]
Thus \(Q_L(\alpha_L)\ge 0\) for all \(\alpha_L\in[0,c]\).

We conclude that
\[
\sigma_H^2 Q_H(\alpha_L)+\sigma_L^2 Q_L(\alpha_L)\ge 0,
\]
and hence
\[
\widetilde W'(\alpha_L)\ge 0
\qquad\text{for all }\alpha_L\in[0,c].
\]

For strict monotonicity, suppose \(n_L>1\) and \(\sigma_L^2>0\), and take
\(\alpha_L\in(0,c)\). Then \(1-\alpha_L>0\), \(t c+\alpha_L>0\), and
\(n_L-1>0\), so the first term in \(Q_L(\alpha_L)\) is strictly positive:
\[
(1-\alpha_L)(t c+\alpha_L)^3(n_L-1)>0.
\]
Hence \(Q_L(\alpha_L)>0\), which implies
\[
\sigma_H^2 Q_H(\alpha_L)+\sigma_L^2 Q_L(\alpha_L)>0.
\]
Therefore
\[
\widetilde W'(\alpha_L)>0
\qquad\text{for all }\alpha_L\in(0,c).
\]
This proves the claim.
\end{proof}

\begin{lemma}
\label{lem:UH-vertical-monotone}
Throughout the feasible region,
\[
\frac{\partial U_H(\alpha_H,\alpha_L)}{\partial \alpha_L}\ge 0.
\]
\end{lemma}

\begin{proof}
Using the expression for \(U_H(\alpha_H,\alpha_L)\) from
[reference to the current equation], write
\[
T=\alpha_H n_H+\alpha_L n_L,
\qquad
P=\alpha_H n_H\mu_H+\alpha_L n_L\mu_L,
\qquad
V=\alpha_H^2 n_H\sigma_H^2+\alpha_L^2 n_L\sigma_L^2 .
\]
Only the terms involving \(T\), \(P\), and \(V\) depend on \(\alpha_L\). Since
\[
\frac{\partial T}{\partial \alpha_L}=n_L,
\qquad
\frac{\partial P}{\partial \alpha_L}=n_L\mu_L,
\qquad
\frac{\partial V}{\partial \alpha_L}=2\alpha_L n_L\sigma_L^2,
\]
differentiation gives
\begin{align*}
\frac{\partial U_H(\alpha_H,\alpha_L)}{\partial \alpha_L}
&=
\frac{n_L\alpha_H}{T^3}
\left[
(\mu_H-\mu_L)n_H\alpha_H T
+
\gamma\alpha_H
\Big(
\sigma_H^2 n_H\alpha_H(2T-\alpha_H n_H)
-\sigma_L^2\alpha_L T
\Big)
\right].
\end{align*}
All terms outside the square brackets are nonnegative. Moreover,
\(\mu_H>\mu_L\), \(\sigma_H^2\ge \sigma_L^2\), and \(T\ge \alpha_H n_H\), so the
bracketed term is weakly nonnegative. Therefore
\[
\frac{\partial U_H(\alpha_H,\alpha_L)}{\partial \alpha_L}\ge 0,
\]
as claimed.
\end{proof}

\section{Feasibility of the full-pooling allocation \((1,1)\)}
\label{app:pooling-feasibility}

The feasibility of \((1,1)\) is determined by the low-risk participation
constraint. Evaluating the IR--\(L\) constraint at \((1,1)\) yields the
equivalent condition
\begin{equation}\label{eq:(1,1)-exact}
    \frac{\mu_H - \mu_L}{\gamma \sigma_L^2}
    \le
    \frac{n_L + n_H}{2 n_H}
    -
    \frac{n_L\sigma_L^2 + n_H\sigma_H^2}
         {2 n_H (n_H+n_L)\sigma_L^2}.
\end{equation}

Since the second term on the right-hand side is strictly positive,
\eqref{eq:(1,1)-exact} is implied by the stronger but simpler sufficient
condition
\begin{equation}\label{eq:(1,1)-simple}
    \frac{\mu_H - \mu_L}{\gamma \sigma_L^2}
    \le
    \frac{n_L + n_H}{2 n_H}.
\end{equation}

Moreover, consider a sequence of economies in which \(n_H,n_L \to \infty\)
with population shares bounded away from zero. Then the correction term
\[
\frac{n_L\sigma_L^2 + n_H\sigma_H^2}
     {2 n_H (n_H+n_L)\sigma_L^2}
\]
is of order \(1/n\) and therefore vanishes asymptotically. Hence, the exact
feasibility condition \eqref{eq:(1,1)-exact} converges to
\[
    \frac{\mu_H - \mu_L}{\gamma \sigma_L^2}
    \le
    \frac{n_L + n_H}{2 n_H}.
\]


\section{Welfare derivatives with respect to \texorpdfstring{\(\alpha_L\)}{alphaL}}
\label{app:alphaL}

This appendix derives the welfare derivative with respect to \(\alpha_L\),
holding \(\alpha_H\) fixed. Its role is twofold. First, it verifies the mean
invariance used in the main text under type-neutral welfare weights. Second, it
provides the derivative formulas needed to study welfare variation along a
frontier parameterized by \(\alpha_L\), including the first-order condition for
an interior optimum and the boundary derivative at \(\alpha_L=0\).

\subsection{Auxiliary objects}
\label{app:alphaL-objects}

Fix \(\alpha_H\) and treat \(\alpha_L\) as the varying argument. Let
\[
T(\alpha_L)=n_H\alpha_H+n_L\alpha_L,
\qquad
T'(\alpha_L)=\frac{\partial T}{\partial \alpha_L}=n_L.
\]

Let
\[
Q(\alpha_L)=\sum_{j=1}^n \alpha_{\tau(j)}X_j
\]
denote the aggregate transferred loss, where \(\tau(j)\in\{H,L\}\) denotes the
type of agent \(j\). Define
\[
M_Q(\alpha_L)=\E[Q],
\qquad
V_Q(\alpha_L)=\Var(Q),
\]
and for each \(k\in\{H,L\}\),
\[
\mu_k=\E[X_k],
\qquad
v_k=\Var(X_k),
\qquad
\Gamma_k(\alpha_L)=\Cov(X_k,Q).
\]

Since only \(\alpha_L\) varies,
\[
M_Q'(\alpha_L)=\frac{\partial M_Q}{\partial \alpha_L}=n_L\mu_L.
\]
Likewise,
\[
\Gamma_k'(\alpha_L)
=
\frac{\partial}{\partial \alpha_L}\Cov(X_k,Q)
=
\sum_{j\in L}\Cov(X_k,X_j),
\]
and
\[
V_Q'(\alpha_L)
=
\frac{\partial}{\partial \alpha_L}\Var(Q)
=
2\sum_{j\in L}\Cov(X_j,Q).
\]

Under type symmetry, these simplify to
\[
\Gamma_H'(\alpha_L)=n_L c_{HL},
\qquad
\Gamma_L'(\alpha_L)=v_L+(n_L-1)c_{LL},
\]
and
\[
V_Q'(\alpha_L)=2n_L\Gamma_L(\alpha_L).
\]

\subsection{Mean invariance}
\label{app:alphaL-mean}

For type \(k\),
\[
\E(\tilde X_k)=(1-\alpha_k)\mu_k+\frac{\alpha_k}{T}M_Q.
\]
Hence
\[
\frac{\partial}{\partial \alpha_L}\left(\frac{M_Q}{T}\right)
=
\frac{M_Q'T-M_QT'}{T^2}
=
\frac{n_L(\mu_LT-M_Q)}{T^2}.
\]
Therefore
\[
\frac{\partial}{\partial \alpha_L}\E(\tilde X_H)
=
\alpha_H\frac{n_L(\mu_LT-M_Q)}{T^2},
\]
while
\[
\frac{\partial}{\partial \alpha_L}\E(\tilde X_L)
=
-\mu_L+\frac{M_Q}{T}
+\alpha_L\frac{n_L(\mu_LT-M_Q)}{T^2}.
\]

Under type-neutral welfare weights, aggregate expected loss is invariant:
\[
-\left(
n_H\frac{\partial}{\partial \alpha_L}\E(\tilde X_H)
+
n_L\frac{\partial}{\partial \alpha_L}\E(\tilde X_L)
\right)=0.
\]
Thus changes in \(\alpha_L\) affect welfare only through the variance terms.
This is the scalar counterpart of budget balance under the sharing rule.

\subsection{Variance derivatives}
\label{app:alphaL-var}

For type \(k\),
\[
\tilde X_k=(1-\alpha_k)X_k+\frac{\alpha_k}{T}Q,
\]
so
\begin{equation}
\Var(\tilde X_k)
=
(1-\alpha_k)^2v_k
+
\left(\frac{\alpha_k}{T}\right)^2V_Q
+
2(1-\alpha_k)\frac{\alpha_k}{T}\Gamma_k.
\label{eq:app-var-general}
\end{equation}

We also use
\[
\frac{\partial}{\partial \alpha_L}\left(\frac{1}{T}\right)
=
-\frac{n_L}{T^2},
\qquad
\frac{\partial}{\partial \alpha_L}\left(\frac{1}{T^2}\right)
=
-\frac{2n_L}{T^3},
\]
\begin{equation}
\frac{\partial}{\partial \alpha_L}\left(\frac{V_Q}{T^2}\right)
=
\frac{V_Q'T-2n_LV_Q}{T^3},
\label{eq:app-D1}
\end{equation}
and
\begin{equation}
\frac{\partial}{\partial \alpha_L}\left(\frac{\Gamma_k}{T}\right)
=
\frac{\Gamma_k'T-n_L\Gamma_k}{T^2}.
\label{eq:app-D2}
\end{equation}

\paragraph{Derivative of \(\Var(\tilde X_H)\).}
Since \(\alpha_H\) is fixed,
\begin{equation}
\frac{\partial}{\partial \alpha_L}\Var(\tilde X_H)
=
\alpha_H^2\frac{V_Q'T-2n_LV_Q}{T^3}
+
2(1-\alpha_H)\alpha_H\frac{\Gamma_H'T-n_L\Gamma_H}{T^2}.
\label{eq:app-dVarH-raw}
\end{equation}

\paragraph{Derivative of \(\Var(\tilde X_L)\).}
Differentiating \eqref{eq:app-var-general} term by term gives
\begin{equation}
\begin{aligned}
\frac{\partial}{\partial \alpha_L}\Var(\tilde X_L)
&=
-2(1-\alpha_L)v_L
+2\alpha_L\frac{V_Q}{T^2}
+\alpha_L^2\frac{V_Q'T-2n_LV_Q}{T^3} \\
&\quad
+2(1-2\alpha_L)\frac{\Gamma_L}{T}
+2(\alpha_L-\alpha_L^2)\frac{\Gamma_L'T-n_L\Gamma_L}{T^2}.
\end{aligned}
\label{eq:app-dVarL-raw}
\end{equation}

Equations \eqref{eq:app-dVarH-raw} and \eqref{eq:app-dVarL-raw} are the general
derivative formulas used in the main text.

\subsection{Type-symmetric covariance expansion}
\label{app:alphaL-typesym}

Under type symmetry,
\[
\Cov(X_i,X_j)=
\begin{cases}
c_{HH}, & i\neq j,\ i,j\in H,\\
c_{LL}, & i\neq j,\ i,j\in L,\\
c_{HL}, & i\in H,\ j\in L.
\end{cases}
\]
Then
\begin{equation}
V_Q
=
\alpha_H^2\bigl[n_Hv_H+n_H(n_H-1)c_{HH}\bigr]
+
\alpha_L^2\bigl[n_Lv_L+n_L(n_L-1)c_{LL}\bigr]
+
2\alpha_H\alpha_Ln_Hn_Lc_{HL},
\label{eq:app-VQ-typesym}
\end{equation}
\begin{equation}
\Gamma_H
=
\alpha_H\bigl[v_H+(n_H-1)c_{HH}\bigr]
+
\alpha_Ln_Lc_{HL},
\label{eq:app-GammaH-typesym}
\end{equation}
and
\begin{equation}
\Gamma_L
=
\alpha_L\bigl[v_L+(n_L-1)c_{LL}\bigr]
+
\alpha_Hn_Hc_{HL}.
\label{eq:app-GammaL-typesym}
\end{equation}
Differentiating yields
\begin{equation}
V_Q'=2n_L\Gamma_L,
\label{eq:app-VQprime-typesym}
\end{equation}
\begin{equation}
\Gamma_H'=n_Lc_{HL},
\qquad
\Gamma_L'=v_L+(n_L-1)c_{LL}.
\label{eq:app-Gammaprime-typesym}
\end{equation}

Substituting
\eqref{eq:app-VQ-typesym}--\eqref{eq:app-Gammaprime-typesym} into
\eqref{eq:app-dVarH-raw} and \eqref{eq:app-dVarL-raw} yields a fully explicit
expression for \(\partial W/\partial\alpha_L\) in terms of
\[
\alpha_H,\ \alpha_L,\ n_H,\ n_L,\ v_H,\ v_L,\ c_{HH},\ c_{LL},\ c_{HL}.
\]
This is the form used for the numerical evaluations discussed in the main text.

\subsection{Derivative of welfare and the first-order condition}
\label{app:alphaL-welfare}

Under type-neutral welfare weights, mean invariance implies that
\begin{equation}
\frac{\partial W}{\partial \alpha_L}
=
-\frac{\gamma}{2}
\left[
n_H\frac{\partial}{\partial \alpha_L}\Var(\tilde X_H)
+
n_L\frac{\partial}{\partial \alpha_L}\Var(\tilde X_L)
\right].
\label{eq:app-Wprime}
\end{equation}
Substituting \eqref{eq:app-dVarH-raw} and \eqref{eq:app-dVarL-raw} into
\eqref{eq:app-Wprime} gives the derivative formula used in
Section~\ref{sec:correlated}.

If welfare is optimized over a frontier parameterized by \(\alpha_L\), then any
interior optimum \(\alpha_L^\star\) must satisfy
\begin{equation}
\frac{\partial W}{\partial \alpha_L}(\alpha_H,\alpha_L^\star)=0,
\label{eq:app-FOC}
\end{equation}
or equivalently,
\begin{equation}
n_H\frac{\partial}{\partial \alpha_L}\Var(\tilde X_H)
+
n_L\frac{\partial}{\partial \alpha_L}\Var(\tilde X_L)
=0.
\label{eq:app-FOC-var}
\end{equation}
Under the type-symmetric covariance parameterization, this condition becomes a
nonlinear equation in \(\alpha_L\), with coefficients determined by the
primitive moments and the fixed value of \(\alpha_H\).

\subsection{Boundary derivative at \texorpdfstring{\(\alpha_L=0\)}{alphaL=0}}
\label{app:alphaL-boundary}

To study whether welfare initially rises as low-risk participation increases
from zero, it is useful to evaluate the derivative at the lower boundary. When
\(\alpha_L=0\),
\[
T=n_H\alpha_H,
\qquad
Q=\alpha_H\sum_{i\in H}X_i,
\]
and under type symmetry,
\[
V_Q
=
\alpha_H^2\bigl[n_Hv_H+n_H(n_H-1)c_{HH}\bigr],
\]
\[
\Gamma_H
=
\alpha_H\bigl[v_H+(n_H-1)c_{HH}\bigr],
\qquad
\Gamma_L
=
\alpha_Hn_Hc_{HL}.
\]

Substituting these expressions into \eqref{eq:app-dVarH-raw} and
\eqref{eq:app-dVarL-raw} yields
\begin{equation}
\frac{\partial}{\partial \alpha_L}\Var(\tilde X_H)\Big|_{\alpha_L=0}
=
-\frac{2\alpha_H}{n_H}
\bigl[v_H+(n_H-1)c_{HH}\bigr]
+
\frac{2(1-\alpha_H)}{n_H}
\Bigl[n_Lc_{HL}-\bigl(v_H+(n_H-1)c_{HH}\bigr)\Bigr],
\label{eq:app-boundary-dVarH}
\end{equation}
and
\begin{equation}
\frac{\partial}{\partial \alpha_L}\Var(\tilde X_L)\Big|_{\alpha_L=0}
=
-2v_L
+
2c_{HL}
+
\frac{2}{n_H}
\bigl[v_H+(n_H-1)c_{HH}\bigr].
\label{eq:app-boundary-dVarL}
\end{equation}
Hence
\begin{equation}
\frac{\partial W}{\partial \alpha_L}\Big|_{\alpha_L=0}
=
-\frac{\gamma}{2}
\left[
n_H\frac{\partial}{\partial \alpha_L}\Var(\tilde X_H)\Big|_{\alpha_L=0}
+
n_L\frac{\partial}{\partial \alpha_L}\Var(\tilde X_L)\Big|_{\alpha_L=0}
\right].
\label{eq:app-boundary-W}
\end{equation}

In contrast to the independent benchmark, the sign of
\(\partial W/\partial\alpha_L\vert_{\alpha_L=0}\) is not generally pinned down
under arbitrary covariance patterns. The expression in
\eqref{eq:app-boundary-W} is therefore best interpreted as a diagnostic object:
it makes transparent how the initial welfare effect of admitting low-risk
agents depends on the trade-off between diversification gains and increased
covariance exposure.



\newpage
\bibliographystyle{apalike}
\bibliography{references}

\end{document}